\documentclass[11pt]{article}

% ---- theme variables (passed via pandoc -V) ----
\newcommand{\ThemeName}{Classic ML Algorithms}
\newcommand{\ThemeKicker}{AI/ML Track}
\newcommand{\ThemeNumber}{02}

\usepackage{interviewstyle}
\setthemecolor{7c3aed}{a78bfa}

% pandoc-required plumbing
\providecommand{\tightlist}{\setlength{\itemsep}{1pt}\setlength{\parskip}{0pt}}
\setlength{\emergencystretch}{3em}
\providecommand{\pandocbounded}[1]{#1}

% syntax-highlighting macros emitted by pandoc (skylighting)
\usepackage{color}
\usepackage{fancyvrb}
\newcommand{\VerbBar}{|}
\newcommand{\VERB}{\Verb[commandchars=\\\{\}]}
\DefineVerbatimEnvironment{Highlighting}{Verbatim}{commandchars=\\\{\}}
% Add ',fontsize=\small' for more characters per line
\usepackage{framed}
\definecolor{shadecolor}{RGB}{35,38,41}
\newenvironment{Shaded}{\begin{snugshade}}{\end{snugshade}}
\newcommand{\AlertTok}[1]{\textcolor[rgb]{0.58,0.85,0.30}{\textbf{\colorbox[rgb]{0.30,0.12,0.14}{#1}}}}
\newcommand{\AnnotationTok}[1]{\textcolor[rgb]{0.25,0.50,0.35}{#1}}
\newcommand{\AttributeTok}[1]{\textcolor[rgb]{0.16,0.50,0.73}{#1}}
\newcommand{\BaseNTok}[1]{\textcolor[rgb]{0.96,0.45,0.00}{#1}}
\newcommand{\BuiltInTok}[1]{\textcolor[rgb]{0.50,0.55,0.55}{#1}}
\newcommand{\CharTok}[1]{\textcolor[rgb]{0.24,0.68,0.91}{#1}}
\newcommand{\CommentTok}[1]{\textcolor[rgb]{0.48,0.49,0.49}{#1}}
\newcommand{\CommentVarTok}[1]{\textcolor[rgb]{0.50,0.55,0.55}{#1}}
\newcommand{\ConstantTok}[1]{\textcolor[rgb]{0.15,0.68,0.68}{\textbf{#1}}}
\newcommand{\ControlFlowTok}[1]{\textcolor[rgb]{0.99,0.74,0.29}{\textbf{#1}}}
\newcommand{\DataTypeTok}[1]{\textcolor[rgb]{0.16,0.50,0.73}{#1}}
\newcommand{\DecValTok}[1]{\textcolor[rgb]{0.96,0.45,0.00}{#1}}
\newcommand{\DocumentationTok}[1]{\textcolor[rgb]{0.64,0.20,0.25}{#1}}
\newcommand{\ErrorTok}[1]{\textcolor[rgb]{0.85,0.27,0.33}{\underline{#1}}}
\newcommand{\ExtensionTok}[1]{\textcolor[rgb]{0.00,0.60,1.00}{\textbf{#1}}}
\newcommand{\FloatTok}[1]{\textcolor[rgb]{0.96,0.45,0.00}{#1}}
\newcommand{\FunctionTok}[1]{\textcolor[rgb]{0.56,0.27,0.68}{#1}}
\newcommand{\ImportTok}[1]{\textcolor[rgb]{0.15,0.68,0.38}{#1}}
\newcommand{\InformationTok}[1]{\textcolor[rgb]{0.77,0.36,0.00}{#1}}
\newcommand{\KeywordTok}[1]{\textcolor[rgb]{0.81,0.81,0.76}{\textbf{#1}}}
\newcommand{\NormalTok}[1]{\textcolor[rgb]{0.81,0.81,0.76}{#1}}
\newcommand{\OperatorTok}[1]{\textcolor[rgb]{0.81,0.81,0.76}{#1}}
\newcommand{\OtherTok}[1]{\textcolor[rgb]{0.15,0.68,0.38}{#1}}
\newcommand{\PreprocessorTok}[1]{\textcolor[rgb]{0.15,0.68,0.38}{#1}}
\newcommand{\RegionMarkerTok}[1]{\textcolor[rgb]{0.16,0.50,0.73}{\colorbox[rgb]{0.08,0.19,0.26}{#1}}}
\newcommand{\SpecialCharTok}[1]{\textcolor[rgb]{0.24,0.68,0.91}{#1}}
\newcommand{\SpecialStringTok}[1]{\textcolor[rgb]{0.85,0.27,0.33}{#1}}
\newcommand{\StringTok}[1]{\textcolor[rgb]{0.96,0.31,0.31}{#1}}
\newcommand{\VariableTok}[1]{\textcolor[rgb]{0.15,0.68,0.68}{#1}}
\newcommand{\VerbatimStringTok}[1]{\textcolor[rgb]{0.85,0.27,0.33}{#1}}
\newcommand{\WarningTok}[1]{\textcolor[rgb]{0.85,0.27,0.33}{#1}}

\begin{document}

\makecoverpage

\thispagestyle{empty}
{\hypersetup{hidelinks}\tableofcontents}
\clearpage

\section{Classic Machine Learning Algorithms}\label{classic-machine-learning-algorithms}

\begin{quote}
Every classic ML algorithm is an answer to the same question: \emph{given past observations, what is the best bet about the future?} --- they differ only in what "best" and "bet" mean.
\end{quote}

\subsection{1. Overview --- What it is}\label{1-overview--what-it-is}

Classic ML algorithms are \textbf{mathematically grounded, interpretable models} that learn patterns from structured/tabular data without requiring deep neural networks. They span supervised learning (regression, classification), unsupervised learning (clustering, dimensionality reduction), and ensemble methods that combine weak learners into strong ones.

\begin{longtable}[]{@{}
  >{\raggedright\arraybackslash}p{(\columnwidth - 2\tabcolsep) * \real{0.3061}}
  >{\raggedright\arraybackslash}p{(\columnwidth - 2\tabcolsep) * \real{0.6939}}@{}}
\toprule\noalign{}
\begin{minipage}[b]{\linewidth}\raggedright
Family
\end{minipage} & \begin{minipage}[b]{\linewidth}\raggedright
Algorithms
\end{minipage} \\
\midrule\noalign{}
\endhead
\bottomrule\noalign{}
\endlastfoot
\textbf{Linear Models} & Linear Regression, Logistic Regression, SVM \\
\textbf{Tree-Based} & Decision Trees, Random Forests, Gradient Boosting (GBM/XGBoost/LightGBM) \\
\textbf{Instance-Based} & k-Nearest Neighbors \\
\textbf{Probabilistic} & Naive Bayes \\
\textbf{Clustering} & k-Means, Hierarchical, DBSCAN \\
\textbf{Dimensionality Reduction} & PCA, t-SNE, UMAP \\
\textbf{Ensembles} & Bagging, Boosting, Stacking \\
\end{longtable}

\textbf{Scope of this document:} First-principles mastery of each family --- the math that drives it, when to reach for it, its failure modes, and exactly how FAANG interviews probe it.

\subsection{2. Why It Exists}\label{2-why-it-exists}

Neural networks solve everything, right? Wrong --- at least in practice:

\begin{itemize}
\tightlist
\item
  \textbf{Tabular data is still king in industry.} Most business data (clickstream, transactions, sensor logs) lives in tables. Tree-based models dominate Kaggle tabular competitions and production ML at FAANG.
\item
  \textbf{Interpretability requirements.} Regulated industries (finance, healthcare) require explainable decisions. Logistic Regression and Decision Trees provide direct feature attribution.
\item
  \textbf{Speed and simplicity.} A Logistic Regression trains in seconds, deploys in microseconds, and never crashes a TPU pod.
\item
  \textbf{Baselines.} Every ML project should start with a linear model baseline. If a neural network can\textquotesingle t beat Ridge Regression on your dataset, you have a data problem, not a model problem.
\item
  \textbf{Feature engineering lever.} Classic models expose \emph{what features matter} --- giving feedback loops that improve data pipelines, which ultimately improve even neural models.
\end{itemize}

The historical arc: \textbf{OLS (1800s) \(\rightarrow\) Logistic Regression (1950s) \(\rightarrow\) Perceptron \(\rightarrow\) SVMs (1990s) \(\rightarrow\) Random Forests (2001) \(\rightarrow\) Gradient Boosting (Friedman 1999, XGBoost 2016) \(\rightarrow\) Transformers (2017)}. Each generation solved a failure mode of the previous one.

\subsection{3. Why FAANG Cares}\label{3-why-faang-cares}

\subsubsection{Google / YouTube}\label{google--youtube}

\begin{itemize}
\tightlist
\item
  \textbf{Gradient Boosting} powers YouTube recommendation ranking (the "scoring" layer after candidate generation). LightGBM scores 100M+ candidates per second in ranking pipelines.
\item
  \textbf{Logistic Regression} with hand-crafted features was the backbone of Google\textquotesingle s ad click-through rate (CTR) prediction for years (FTRL-Proximal algorithm, Google 2013 paper).
\end{itemize}

\subsubsection{Meta}\label{meta}

\begin{itemize}
\tightlist
\item
  \textbf{GBDT + Neural Net hybrid (GBDT-NN):} Meta\textquotesingle s ad ranking uses gradient boosted trees to generate transformed features that feed into a neural network --- the trees capture non-linear feature interactions, the net captures higher-order patterns.
\item
  \textbf{k-Means:} User clustering for audience segmentation and look-alike modeling.
\end{itemize}

\subsubsection{Amazon}\label{amazon}

\begin{itemize}
\tightlist
\item
  \textbf{Random Forests} in fraud detection (interpretable feature importance for chargebacks).
\item
  \textbf{k-NN} in product recommendation ("customers who bought X also bought Y" --- literally nearest neighbors in embedding space).
\end{itemize}

\subsubsection{Apple}\label{apple}

\begin{itemize}
\tightlist
\item
  \textbf{Naive Bayes} still used in spam filtering (iCloud Mail). Fast, privacy-preserving (can train on-device).
\item
  \textbf{PCA} for face ID feature compression before downstream classifiers.
\end{itemize}

\subsubsection{Microsoft (Azure ML / LinkedIn)}\label{microsoft-azure-ml--linkedin}

\begin{itemize}
\tightlist
\item
  \textbf{SVM} historically used in NLP classification tasks (text SVMs with TF-IDF features).
\item
  \textbf{LightGBM} is a Microsoft Research product --- used across Bing ranking.
\end{itemize}

\textbf{Key interview insight:} FAANG doesn\textquotesingle t ask "which algorithm is best?" --- they ask "which algorithm is right \emph{for this constraint set}?" Knowing where each lives in production is what separates senior candidates.

\subsection{4. Core Concepts}\label{4-core-concepts}

\subsubsection{4.1 Linear Regression}\label{41-linear-regression}

\paragraph{Intuition}\label{intuition}

Fit the \textbf{best straight line} (hyperplane in higher dimensions) through data points to predict a continuous output. "Best" = minimizing squared residuals.

\paragraph{The Math}\label{the-math}

\textbf{Model:} \texttt{y\_hat\ =\ w\^{}T\ x\ +\ b} where \texttt{w} = weights, \texttt{x} = features, \texttt{b} = bias

\textbf{Loss (OLS --- Ordinary Least Squares):}

\begin{verbatim}
L(w) = (1/n) * sum[(y_i - w^T x_i)^2]
     = (1/n) * ||y - Xw||^2
\end{verbatim}

\textbf{Closed-form solution (Normal Equation):}

\begin{verbatim}
w* = (X^T X)^{-1} X^T y
\end{verbatim}

\begin{itemize}
\tightlist
\item
  Works when \texttt{X\^{}T\ X} is invertible (no perfectly collinear features).
\item
  Time complexity: O(n * p\^{}2 + p\^{}3) --- cubic in features, so use gradient descent when p \textgreater{} \textasciitilde10,000.
\end{itemize}

\textbf{Gradient descent update:}

\begin{verbatim}
w := w - alpha * (2/n) * X^T(Xw - y)
\end{verbatim}

\paragraph{Gauss-Markov Assumptions (OLS is BLUE --- Best Linear Unbiased Estimator)}\label{gauss-markov-assumptions-ols-is-blue--best-linear-unbiased-estimator}

\begin{enumerate}
\def\labelenumi{\arabic{enumi}.}
\tightlist
\item
  \textbf{Linearity:} \texttt{y\ =\ X*beta\ +\ epsilon}
\item
  \textbf{No multicollinearity:} features are not perfectly linearly dependent
\item
  \textbf{Homoscedasticity:} \texttt{Var(epsilon\_i)\ =\ sigma\^{}2} (constant variance)
\item
  \textbf{No autocorrelation:} \texttt{Cov(epsilon\_i,\ epsilon\_j)\ =\ 0} for \texttt{i\ !=\ j}
\item
  \textbf{Exogeneity:} \texttt{E{[}epsilon\ \textbar{}\ X{]}\ =\ 0} (errors uncorrelated with features)
\item
  \emph{(Optional for inference)} Normality: \texttt{epsilon\ \textasciitilde{}\ N(0,\ sigma\^{}2)}
\end{enumerate}

\textbf{Violation consequences:}
\textbar{} Violation \textbar{} Consequence \textbar{} Fix \textbar{}
\textbar-\/-\/-\textbar-\/-\/-\textbar-\/-\/-\textbar{}
\textbar{} Multicollinearity \textbar{} High variance in w estimates, wrong signs \textbar{} Ridge/Lasso regularization \textbar{}
\textbar{} Heteroscedasticity \textbar{} Inefficient estimates, wrong standard errors \textbar{} Weighted LS, log-transform y \textbar{}
\textbar{} Non-linearity \textbar{} Biased predictions \textbar{} Polynomial features, tree models \textbar{}
\textbar{} Autocorrelation \textbar{} Underestimated variance \textbar{} Time-series models (ARIMA), GLS \textbar{}

\paragraph{Regularization}\label{regularization}

\begin{verbatim}
Ridge (L2):  L = ||y - Xw||^2 + lambda * ||w||^2       (shrinks all weights)
Lasso (L1):  L = ||y - Xw||^2 + lambda * ||w||_1       (sparsifies weights)
ElasticNet:  L = ||y - Xw||^2 + lambda1*||w||_1 + lambda2*||w||^2
\end{verbatim}

\textbf{Ridge} has closed form: \texttt{w*\ =\ (X\^{}T\ X\ +\ lambda*I)\^{}\{-1\}\ X\^{}T\ y} --- the \texttt{lambda*I} term fixes invertibility even under multicollinearity.

\textbf{Lasso} encourages sparsity because the L1 ball has corners at the axes --- the constrained solution tends to land at corners (where some weights = 0). This is automatic feature selection.

\textbf{Key hyperparameters:} \texttt{alpha} (learning rate), \texttt{lambda} (regularization strength), number of iterations

\textbf{Complexity:} Training O(n\emph{p\^{}2 + p\^{}3) closed form, O(n}p*iter) gradient descent; Inference O(p)

\textbf{Strengths:} Fast, interpretable, works on very high-dimensional sparse data (with regularization), provides uncertainty estimates (via confidence intervals)

\textbf{Weaknesses:} Assumes linearity, sensitive to outliers (squared loss), needs feature scaling for gradient descent

\subsubsection{4.2 Logistic Regression}\label{42-logistic-regression}

\paragraph{Intuition}\label{intuition-1}

Model the \textbf{probability} that an input belongs to class 1. Use a linear model internally, then squeeze the output into {[}0,1{]} with the sigmoid function.

\paragraph{The Math}\label{the-math-1}

\textbf{Sigmoid:}

\begin{verbatim}
sigma(z) = 1 / (1 + e^{-z})

Properties:
- sigma(0) = 0.5
- sigma(+inf) = 1, sigma(-inf) = 0
- d/dz sigma(z) = sigma(z) * (1 - sigma(z))   <-- elegant gradient
\end{verbatim}

\textbf{Model:}

\begin{verbatim}
P(y=1 | x) = sigma(w^T x + b)
\end{verbatim}

\textbf{Log-Loss (Binary Cross-Entropy):}

\begin{verbatim}
L(w) = -(1/n) * sum[ y_i * log(p_i) + (1-y_i) * log(1-p_i) ]
\end{verbatim}

Why log-loss? It\textquotesingle s the \textbf{negative log-likelihood} under a Bernoulli distribution. Maximizing likelihood = minimizing log-loss.

\textbf{No closed form} --- must use gradient descent:

\begin{verbatim}
dL/dw = (1/n) * X^T (p - y)
        where p_i = sigma(w^T x_i)
\end{verbatim}

The gradient has the same form as linear regression --- clean and memorable.

\paragraph{Decision Boundary}\label{decision-boundary}

\begin{verbatim}
P(y=1|x) >= 0.5  when  w^T x + b >= 0
\end{verbatim}

The boundary is the hyperplane \texttt{w\^{}T\ x\ +\ b\ =\ 0} --- \textbf{linear} in feature space. For non-linear boundaries, add polynomial features or use kernels.

\paragraph{Interpretation: Log-Odds (Logit)}\label{interpretation-log-odds-logit}

\begin{verbatim}
log[P(y=1)/(1-P(y=1))] = w^T x + b
\end{verbatim}

Each weight \texttt{w\_j} tells you: "a unit increase in feature \texttt{x\_j} multiplies the \textbf{odds} by \texttt{e\^{}\{w\_j\}}." This is why logistic regression is the workhorse in medicine and social science --- coefficients have direct interpretations.

\paragraph{Multi-class: Softmax Regression}\label{multi-class-softmax-regression}

\begin{verbatim}
P(y=k | x) = e^{w_k^T x} / sum_j[e^{w_j^T x}]
Loss: Cross-Entropy = -(1/n) * sum_i sum_k [y_{ik} * log(p_{ik})]
\end{verbatim}

\textbf{Key hyperparameters:} Regularization (C = 1/lambda in sklearn), solver (lbfgs, liblinear, saga), max\_iter, class\_weight

\textbf{Strengths:} Probabilistic output, fast, interpretable weights, works on high-dim sparse data (text), no feature scaling required (but helps convergence)

\textbf{Weaknesses:} Assumes linear decision boundary, poor with feature interactions unless manually added, sensitive to outlier features

\textbf{When to use:} Binary/multiclass classification, when you need calibrated probabilities, as a baseline, in production when inference speed matters, when interpretability is legally required

\subsubsection{4.3 Decision Trees}\label{43-decision-trees}

\paragraph{Intuition}\label{intuition-2}

Recursively \textbf{split} the data using the feature that best separates classes/reduces prediction error. Creates a tree where each internal node is a test, each leaf is a prediction.

\paragraph{Splitting Criteria}\label{splitting-criteria}

\textbf{For Classification:}

\textbf{Entropy (Information Gain):}

\begin{verbatim}
Entropy(S) = -sum_k [ p_k * log2(p_k) ]

Information Gain(S, A) = Entropy(S) - sum_{v in vals(A)} [ |S_v|/|S| * Entropy(S_v) ]
\end{verbatim}

\begin{itemize}
\tightlist
\item
  Entropy = 0 when all samples are same class (pure)
\item
  Entropy = 1 for binary with 50/50 split (maximum impurity)
\end{itemize}

\textbf{Gini Impurity:}

\begin{verbatim}
Gini(S) = 1 - sum_k [ p_k^2 ]
         = sum_k [ p_k * (1 - p_k) ]
\end{verbatim}

\begin{itemize}
\tightlist
\item
  Gini = 0 when pure, Gini = 0.5 for 50/50 binary split
\item
  \textbf{Gini is computationally cheaper} (no log), used in sklearn by default
\item
  \textbf{Entropy is more sensitive} to class probability changes --- theoretically better for imbalanced problems
\end{itemize}

\textbf{Mathematically:} Gini approximates entropy via \texttt{ln(x)\ ≈\ (x-1)}. In practice, results are nearly identical.

\textbf{For Regression:}

\begin{verbatim}
MSE split: choose feature/threshold minimizing sum of within-child MSE
Variance reduction: Var(parent) - weighted_avg[Var(children)]
\end{verbatim}

\paragraph{The Splitting Algorithm (CART)}\label{the-splitting-algorithm-cart}

\begin{verbatim}
function BuildTree(data, depth):
    if stopping_criterion_met:
        return Leaf(majority_class or mean(y))
    
    best_feature, best_threshold = None, None
    best_gain = -inf
    
    for each feature f:
        for each candidate threshold t:
            gain = impurity(data) - weighted_impurity(split(data, f, t))
            if gain > best_gain:
                best_gain = gain
                best_feature, best_threshold = f, t
    
    left, right = split(data, best_feature, best_threshold)
    return Node(best_feature, best_threshold,
                left=BuildTree(left, depth+1),
                right=BuildTree(right, depth+1))
\end{verbatim}

\textbf{Candidate thresholds:} Sort feature values, try midpoints between consecutive unique values. O(n log n) per feature.

\textbf{Total training complexity:} O(n * p * log(n) * depth) per node, roughly O(n * p * log\^{}2(n)) total

\paragraph{Pruning}\label{pruning}

\textbf{Pre-pruning (early stopping):}

\begin{itemize}
\tightlist
\item
  \texttt{max\_depth}: limit tree depth
\item
  \texttt{min\_samples\_split}: minimum samples to split a node
\item
  \texttt{min\_samples\_leaf}: minimum samples in a leaf
\item
  \texttt{min\_impurity\_decrease}: only split if gain exceeds threshold
\end{itemize}

\textbf{Post-pruning (cost-complexity pruning):}

\begin{verbatim}
Minimize: sum_leaves[impurity(leaf)] + alpha * num_leaves
\end{verbatim}

\begin{itemize}
\tightlist
\item
  Grow full tree, then prune leaves bottom-up
\item
  \texttt{alpha} (ccp\_alpha in sklearn) controls complexity penalty
\item
  Cross-validate to find best alpha
\end{itemize}

\paragraph{Strengths}\label{strengths}

\begin{itemize}
\tightlist
\item
  Interpretable (can print the tree, generate rules)
\item
  Handles mixed feature types naturally
\item
  Handles missing values (surrogate splits)
\item
  No feature scaling needed
\item
  Captures non-linear interactions
\end{itemize}

\paragraph{Weaknesses}\label{weaknesses}

\begin{itemize}
\tightlist
\item
  \textbf{High variance} --- small data changes create very different trees (instability)
\item
  \textbf{Greedy} --- local splits are not globally optimal
\item
  Tends to overfit without pruning
\item
  Biased toward high-cardinality features (Gini handles this better than IG)
\item
  \textbf{Piecewise constant prediction} --- bad at smooth regression
\end{itemize}

\textbf{Key hyperparameters:} max\_depth, min\_samples\_split, min\_samples\_leaf, criterion (gini/entropy), ccp\_alpha

\subsubsection{4.4 Random Forests}\label{44-random-forests}

\paragraph{Intuition}\label{intuition-3}

Train \textbf{many decision trees on random subsets of data and features}, then aggregate their predictions. The randomness decorrelates the trees --- their errors cancel out.

\paragraph{The Algorithm: Bagging + Feature Subsampling}\label{the-algorithm-bagging--feature-subsampling}

\begin{verbatim}
function RandomForest(X, y, n_trees, max_features):
    trees = []
    for i in 1..n_trees:
        # Bootstrap sample
        X_boot, y_boot = bootstrap(X, y)          # sample n rows WITH replacement
        
        # Build tree with feature subsampling
        tree = BuildTree(X_boot, y_boot,
                         max_features=sqrt(p))     # at each split, consider only sqrt(p) features
        trees.append(tree)
    
    return trees

function Predict(x, trees):
    predictions = [tree.predict(x) for tree in trees]
    return majority_vote(predictions)              # or mean() for regression
\end{verbatim}

\paragraph{Why Does This Work?}\label{why-does-this-work}

\textbf{Bias-Variance Tradeoff:}

\begin{verbatim}
E[(y - y_hat)^2] = Bias^2 + Variance + Irreducible_Noise

Single deep tree:  Low Bias,  HIGH Variance  (overfits)
Random Forest:     Low Bias,  Low Variance   (errors average out)
\end{verbatim}

\textbf{Mathematical intuition:} If we have \texttt{B} trees each with variance \texttt{sigma\^{}2} and pairwise correlation \texttt{rho}:

\begin{verbatim}
Var(average) = rho * sigma^2 + (1-rho)/B * sigma^2
\end{verbatim}

\begin{itemize}
\tightlist
\item
  If trees were identical (rho=1): no variance reduction
\item
  If trees were independent (rho=0): variance reduces by 1/B
\item
  Random feature subsampling \textbf{reduces rho} without increasing bias much
\end{itemize}

\paragraph{Out-Of-Bag (OOB) Error}\label{out-of-bag-oob-error}

Each bootstrap sample leaves out \textasciitilde36.8\% of data (\texttt{(1-1/n)\^{}n\ →\ 1/e}). These OOB samples form a free validation set:

\begin{itemize}
\tightlist
\item
  OOB error \(\approx\) test error without needing a held-out set
\item
  Useful when data is scarce
\end{itemize}

\paragraph{Feature Importance}\label{feature-importance}

\begin{verbatim}
Importance(feature_j) = sum_trees sum_nodes_using_j [
    weighted_impurity_decrease(node)  # weighted by n_samples
] / n_trees
\end{verbatim}

\textbf{Mean Decrease in Impurity (MDI)} --- fast but biased toward high-cardinality features.
\textbf{Permutation Importance} --- shuffle feature j in validation set, measure accuracy drop. Slower but unbiased.

\textbf{Key hyperparameters:} n\_estimators, max\_features (sqrt(p) for classification, p/3 for regression), max\_depth, min\_samples\_leaf, bootstrap, oob\_score

\textbf{Complexity:} Training O(B * n * sqrt(p) * log\^{}2(n)), Inference O(B * depth)

\textbf{Strengths:} Robust to overfitting, handles high-dim data, feature importance, OOB error, parallel training, handles missing values reasonably, works well out-of-the-box

\textbf{Weaknesses:} Less interpretable than single tree, slow inference for very large forests, memory-intensive, can overfit on noisy datasets, not great at extrapolation

\subsubsection{4.5 Gradient Boosting (GBM / XGBoost / LightGBM)}\label{45-gradient-boosting-gbm--xgboost--lightgbm}

\paragraph{Intuition}\label{intuition-4}

Train trees \textbf{sequentially}, where each tree corrects the errors (residuals) of the previous ensemble. Gradient boosting generalizes this: each tree fits the \textbf{negative gradient} of the loss function.

\paragraph{The Core Algorithm}\label{the-core-algorithm}

\begin{verbatim}
Initialize: F_0(x) = argmin_gamma sum_i L(y_i, gamma)
            (e.g., mean(y) for MSE, log(p/(1-p)) for log-loss)

For m = 1 to M:
    1. Compute pseudo-residuals (negative gradient):
       r_{im} = -[dL(y_i, F(x_i)) / dF(x_i)]    for all i
       
       For MSE: r_{im} = y_i - F_{m-1}(x_i)      (just the residuals!)
       For log-loss: r_{im} = y_i - sigma(F_{m-1}(x_i))
    
    2. Fit a regression tree h_m(x) to {(x_i, r_{im})}
    
    3. Find optimal step size:
       gamma_m = argmin_gamma sum_i L(y_i, F_{m-1}(x_i) + gamma * h_m(x_i))
    
    4. Update:
       F_m(x) = F_{m-1}(x) + learning_rate * gamma_m * h_m(x)

Final: F_M(x)
\end{verbatim}

\textbf{Key insight:} We\textquotesingle re doing \textbf{gradient descent in function space} --- instead of updating weights, we\textquotesingle re updating the entire function by adding a new tree in the direction that reduces loss.

\paragraph{Why Gradient, Not Just Residuals?}\label{why-gradient-not-just-residuals}

For MSE loss, \texttt{dL/dF\ =\ -(y\ -\ F(x))} --- the negative gradient IS the residuals. For other losses (MAE, log-loss, Huber), the "residuals" change, but the gradient descent framework handles them uniformly.

\paragraph{XGBoost Improvements Over Vanilla GBM}\label{xgboost-improvements-over-vanilla-gbm}

\textbf{1. Regularized objective:}

\begin{verbatim}
Obj = sum_i L(y_i, y_hat_i) + sum_k Omega(f_k)
Omega(f) = gamma*T + (1/2)*lambda*||w||^2
\end{verbatim}

Where T = number of leaves, w = leaf weights. This directly penalizes tree complexity.

\textbf{2. Approximate greedy algorithm:}
Instead of trying all split points, use quantile sketches to find split candidates. Much faster for large datasets.

\textbf{3. Sparsity-aware split finding:}
Learn default direction for missing values --- handles sparse features efficiently.

\textbf{4. Column block for parallel learning:}
Pre-sort features, store in compressed column blocks. Enables parallel split finding across features.

\textbf{5. Cache-aware access and out-of-core computation:}
Designed for data that doesn\textquotesingle t fit in RAM.

\paragraph{LightGBM Improvements Over XGBoost}\label{lightgbm-improvements-over-xgboost}

\textbf{1. Gradient-based One-Side Sampling (GOSS):}
Keep all large-gradient samples (they have high information), randomly downsample small-gradient samples. Faster with minimal accuracy loss.

\textbf{2. Exclusive Feature Bundling (EFB):}
Bundle mutually exclusive sparse features into one dense feature. Reduces effective number of features.

\textbf{3. Leaf-wise tree growth (vs level-wise in XGBoost):}

\begin{verbatim}
Level-wise (XGBoost):        Leaf-wise (LightGBM):
Level 0:  [root]             Always split the leaf with max delta_loss
Level 1:  [L] [R]            Can create deep, asymmetric trees
Level 2:  [LL][LR][RL][RR]   More accurate but can overfit
\end{verbatim}

Leaf-wise is faster (fewer splits to get same accuracy) but needs \texttt{num\_leaves} control.

\textbf{4. Histogram-based splitting:}
Bin continuous features into 255 buckets. Huge memory and speed win.

\paragraph{Comparison Table}\label{comparison-table}

\begin{longtable}[]{@{}
  >{\raggedright\arraybackslash}p{(\columnwidth - 8\tabcolsep) * \real{0.1961}}
  >{\raggedright\arraybackslash}p{(\columnwidth - 8\tabcolsep) * \real{0.1471}}
  >{\raggedright\arraybackslash}p{(\columnwidth - 8\tabcolsep) * \real{0.2255}}
  >{\raggedright\arraybackslash}p{(\columnwidth - 8\tabcolsep) * \real{0.1569}}
  >{\raggedright\arraybackslash}p{(\columnwidth - 8\tabcolsep) * \real{0.2745}}@{}}
\toprule\noalign{}
\begin{minipage}[b]{\linewidth}\raggedright
Property
\end{minipage} & \begin{minipage}[b]{\linewidth}\raggedright
GBM
\end{minipage} & \begin{minipage}[b]{\linewidth}\raggedright
XGBoost
\end{minipage} & \begin{minipage}[b]{\linewidth}\raggedright
LightGBM
\end{minipage} & \begin{minipage}[b]{\linewidth}\raggedright
CatBoost
\end{minipage} \\
\midrule\noalign{}
\endhead
\bottomrule\noalign{}
\endlastfoot
\textbf{Speed} & Slow & Medium & Fast & Medium \\
\textbf{Memory} & High & Medium & Low & Medium \\
\textbf{GPU support} & No & Yes & Yes & Yes \\
\textbf{Categorical features} & Manual encoding & Manual encoding & Native (limited) & Native (excellent) \\
\textbf{Tree growth} & Level-wise & Level-wise & Leaf-wise & Symmetric \\
\textbf{Regularization} & Shrinkage & L1, L2, tree complexity & L1, L2 & Leaf-wise + ordered boosting \\
\textbf{Missing values} & Manual & Native & Native & Native \\
\textbf{Best for} & Baseline & General & Large datasets & Categorical-heavy \\
\end{longtable}

\textbf{Key hyperparameters:}

\begin{itemize}
\tightlist
\item
  \texttt{n\_estimators} (M): number of trees --- more = less bias, more variance, slower
\item
  \texttt{learning\_rate} (eta): step size --- smaller = more robust, needs more trees
\item
  \texttt{max\_depth}: tree depth --- shallower = less variance, more bias
\item
  \texttt{subsample}: row subsampling per tree (stochastic GBM) --- reduces variance
\item
  \texttt{colsample\_bytree}: column subsampling --- like RF feature randomness
\item
  \texttt{min\_child\_weight} / \texttt{min\_child\_samples}: regularization on leaves
\item
  \texttt{lambda/alpha}: L2/L1 regularization on leaf weights
\end{itemize}

\textbf{Golden rule:} \texttt{learning\_rate\ *\ n\_estimators} should be roughly constant. Lower LR + more trees = better but slower.

\textbf{Complexity:} Training O(M * n * p * log(n)), Inference O(M * depth)

\textbf{Strengths:} State-of-the-art on tabular data, handles heterogeneous features, built-in regularization, feature importance, handles missing values, early stopping

\textbf{Weaknesses:} Many hyperparameters, sequential training (harder to parallelize than RF), can overfit with wrong params, slow to train at scale, requires careful tuning

\subsubsection{4.6 Support Vector Machines (SVM)}\label{46-support-vector-machines-svm}

\paragraph{Intuition}\label{intuition-5}

Find the \textbf{hyperplane that maximizes the margin} between classes. The margin is the distance from the hyperplane to the nearest training point of each class. Wider margin = better generalization.

\paragraph{Hard-Margin SVM (linearly separable)}\label{hard-margin-svm-linearly-separable}

\textbf{Objective:}

\begin{verbatim}
Maximize margin = 2 / ||w||
Subject to: y_i(w^T x_i + b) >= 1   for all i
Equivalently: Minimize (1/2)||w||^2
              Subject to: y_i(w^T x_i + b) >= 1
\end{verbatim}

The hyperplane is \texttt{w\^{}T\ x\ +\ b\ =\ 0}.\\
The two margin planes are \texttt{w\^{}T\ x\ +\ b\ =\ +1} and \texttt{w\^{}T\ x\ +\ b\ =\ -1}.\\
Margin width = \texttt{2/\textbar{}\textbar{}w\textbar{}\textbar{}}.

\textbf{Support Vectors:} The training points that lie exactly on the margin planes (\texttt{y\_i(w\^{}T\ x\_i\ +\ b)\ =\ 1}). Only these points determine the hyperplane --- the rest are irrelevant. This is why SVMs are robust to outliers far from the boundary.

\paragraph{Soft-Margin SVM (C-SVM, non-separable)}\label{soft-margin-svm-c-svm-non-separable}

Allow some points to violate the margin via \textbf{slack variables} \texttt{xi\_i\ \textgreater{}=\ 0}:

\begin{verbatim}
Minimize: (1/2)||w||^2 + C * sum_i xi_i
Subject to: y_i(w^T x_i + b) >= 1 - xi_i
            xi_i >= 0
\end{verbatim}

\textbf{C = regularization parameter:}

\begin{itemize}
\tightlist
\item
  \texttt{C\ →\ ∞}: Hard margin, no violations allowed, risk of overfitting
\item
  \texttt{C\ →\ 0}: Wide margin, many violations allowed, underfitting
\item
  \textbf{Smaller C = more regularization} (opposite intuition from other models)
\end{itemize}

\textbf{Hinge Loss equivalence:}

\begin{verbatim}
Hinge Loss: max(0, 1 - y_i * f(x_i))
SVM objective = (1/2)||w||^2 + C * sum_i max(0, 1 - y_i * f(x_i))
             = L2 regularization + C * sum of hinge losses
\end{verbatim}

\paragraph{The Kernel Trick}\label{the-kernel-trick}

Map data to higher dimensions where it becomes linearly separable --- without ever computing the high-dimensional vectors.

\textbf{Key insight:} The SVM dual formulation depends on data only through \textbf{inner products} \texttt{x\_i\^{}T\ x\_j}:

\begin{verbatim}
Dual: Maximize sum_i alpha_i - (1/2) sum_{ij} alpha_i alpha_j y_i y_j (x_i^T x_j)
\end{verbatim}

Replace \texttt{x\_i\^{}T\ x\_j} with \texttt{K(x\_i,\ x\_j)} --- a \textbf{kernel function} that computes the inner product in some high/infinite-dimensional space, without going there explicitly.

\textbf{Common kernels:}

\begin{longtable}[]{@{}
  >{\raggedright\arraybackslash}p{(\columnwidth - 4\tabcolsep) * \real{0.1424}}
  >{\raggedright\arraybackslash}p{(\columnwidth - 4\tabcolsep) * \real{0.3561}}
  >{\raggedright\arraybackslash}p{(\columnwidth - 4\tabcolsep) * \real{0.5015}}@{}}
\toprule\noalign{}
\begin{minipage}[b]{\linewidth}\raggedright
Kernel
\end{minipage} & \begin{minipage}[b]{\linewidth}\raggedright
Formula
\end{minipage} & \begin{minipage}[b]{\linewidth}\raggedright
When to use
\end{minipage} \\
\midrule\noalign{}
\endhead
\bottomrule\noalign{}
\endlastfoot
\textbf{Linear} & \texttt{K(x,z)\ =\ x\^{}T\ z} & Linearly separable, high-dim sparse (text) \\
\textbf{Polynomial} & \texttt{K(x,z)\ =\ (gamma*x\^{}T\ z\ +\ r)\^{}d} & Moderate non-linearity, image features \\
\textbf{RBF/Gaussian} & `K(x,z) = exp(-gamma * & \\
\textbf{Sigmoid} & \texttt{K(x,z)\ =\ tanh(gamma*x\^{}T\ z\ +\ r)} & Similar to neural net activation, rarely used \\
\end{longtable}

\textbf{RBF Intuition:} \texttt{K(x,z)} measures similarity --- it\textquotesingle s 1 when x=z, 0 when they\textquotesingle re infinitely far apart. Points "vote" on classification weighted by similarity. The implicit feature space is \textbf{infinite-dimensional}.

\textbf{The Kernel Trick computation:}

\begin{verbatim}
Normal: compute phi(x) [infinite dim vector], then dot product --> impossible
Kernel: compute K(x,z) = <phi(x), phi(z)> directly --> O(n_features), no expansion needed
\end{verbatim}

\textbf{Key hyperparameters:} C, kernel type, gamma (for RBF: smaller gamma = wider RBF, smoother boundary), degree (polynomial), coef0

\textbf{Complexity:} Training O(n\^{}2 * p) to O(n\^{}3 * p), Inference O(n\_sv * p) --- scales poorly with large n

\textbf{Strengths:} Works in high-dim spaces, effective when n \textless{} p, memory efficient (only support vectors stored), kernel trick for non-linear data, global optimum (convex problem)

\textbf{Weaknesses:} Slow on large datasets (O(n\^{}2-3)), no probability output (need Platt scaling), sensitive to feature scaling, hard to interpret, tricky to tune gamma and C

\textbf{When to use:} Small-to-medium dataset, many features (text), non-linear classification with RBF, when you need a clear margin

\subsubsection{4.7 k-Nearest Neighbors (k-NN)}\label{47-k-nearest-neighbors-k-nn}

\paragraph{Intuition}\label{intuition-6}

To classify a new point, find its \textbf{k nearest neighbors} in the training set and take a majority vote (classification) or average (regression).

\textbf{The simplest possible model --- no training, just memorization.}

\paragraph{The Algorithm}\label{the-algorithm}

\begin{verbatim}
Training: Store all (x_i, y_i)            # O(n*p) space

Prediction for query q:
1. Compute distance(q, x_i) for all i     # O(n*p)
2. Sort by distance                        # O(n log n)
3. Take k smallest                         # O(k)
4. Return majority vote or mean            # O(k)
\end{verbatim}

\textbf{Distance metrics:}

\begin{verbatim}
Euclidean: d(x,z) = sqrt(sum_j (x_j - z_j)^2)   -- most common
Manhattan: d(x,z) = sum_j |x_j - z_j|            -- robust to outliers
Minkowski: d(x,z) = (sum_j |x_j - z_j|^p)^{1/p} -- generalizes both
Cosine:    d(x,z) = 1 - (x^T z)/(||x||*||z||)   -- text/embedding similarity
Hamming:   for binary/categorical features
\end{verbatim}

\paragraph{Choosing k}\label{choosing-k}

\begin{verbatim}
k=1:       Zero training error, very high variance (memorizes noise)
k=n:       Predict majority class always -- zero variance, high bias
k=sqrt(n): Common heuristic, balance point
\end{verbatim}

\textbf{Small k = high variance (overfitting), Large k = high bias (underfitting)}

Use cross-validation to choose k. Odd k avoids ties in binary classification.

\paragraph{Curse of Dimensionality}\label{curse-of-dimensionality}

In high dimensions, \textbf{all points become equidistant} --- nearest neighbors are no more similar than random points. The ratio of max to min distance converges to 1 as dimensionality grows.

\begin{verbatim}
d_max / d_min → 1 as p → ∞
\end{verbatim}

Consequence: k-NN degrades badly in high dimensions. Always apply PCA/feature selection first.

\paragraph{Improvements}\label{improvements}

\begin{itemize}
\tightlist
\item
  \textbf{KD-tree:} Organize data in a tree for O(p log n) query instead of O(np). Breaks down for p \textgreater{} \textasciitilde20.
\item
  \textbf{Ball-tree:} Better than KD-tree for high dimensions and non-Euclidean metrics.
\item
  \textbf{Approximate Nearest Neighbors (ANN):} FAISS, HNSW, Annoy --- trade exact for speed at scale.
\end{itemize}

\textbf{Key hyperparameters:} k, distance metric, weights (uniform vs distance-weighted), algorithm (ball\_tree, kd\_tree, brute)

\textbf{Complexity:} Training O(n\emph{p) (just storing), Inference O(n}p + n log n) per query

\textbf{Strengths:} Simple, interpretable, no training, naturally handles multi-class, works well with enough data

\textbf{Weaknesses:} Slow at inference, memory-intensive (stores all data), sensitive to irrelevant features, needs feature scaling, curse of dimensionality

\textbf{When to use:} Small datasets, recommendation systems (in embedding space), anomaly detection, when decision boundary is complex and data is low-dimensional

\subsubsection{4.8 Naive Bayes}\label{48-naive-bayes}

\paragraph{Intuition}\label{intuition-7}

Apply Bayes\textquotesingle{} theorem to classify, with the \textbf{"naive" assumption that features are conditionally independent} given the class. Despite this usually being wrong, it works surprisingly well.

\paragraph{Bayes\textquotesingle{} Theorem}\label{bayes-theorem}

\begin{verbatim}
P(y=k | x) = P(x | y=k) * P(y=k) / P(x)
           ∝ P(x | y=k) * P(y=k)          # P(x) is constant for all classes
\end{verbatim}

\textbf{The naive assumption:}

\begin{verbatim}
P(x | y=k) = prod_j P(x_j | y=k)         # features independent given class
\end{verbatim}

\textbf{Classification rule:}

\begin{verbatim}
y_hat = argmax_k [log P(y=k) + sum_j log P(x_j | y=k)]
\end{verbatim}

Log-sum is numerically stable and avoids underflow from multiplying many small probabilities.

\paragraph{Variants}\label{variants}

\textbf{Gaussian NB (continuous features):}

\begin{verbatim}
P(x_j | y=k) = N(x_j; mu_{jk}, sigma_{jk}^2)
mu_{jk} = mean of feature j in class k
sigma_{jk}^2 = variance of feature j in class k
\end{verbatim}

\textbf{Multinomial NB (count features, text):}

\begin{verbatim}
P(x_j | y=k) = theta_{jk}^{x_j}      # x_j is count of word j
theta_{jk} = (count of word j in class k + alpha) / (total words in class k + alpha * |V|)
\end{verbatim}

\texttt{alpha} is Laplace/additive smoothing --- prevents zero probabilities for unseen words.

\textbf{Bernoulli NB (binary features):}

\begin{verbatim}
P(x_j | y=k) = p_{jk}^{x_j} * (1-p_{jk})^{1-x_j}
\end{verbatim}

\paragraph{Why It Works Despite Wrong Assumption}\label{why-it-works-despite-wrong-assumption}

\begin{itemize}
\tightlist
\item
  For classification, we only need \textbf{ordering} of P(y=k\textbar x) across classes, not exact values
\item
  Independence errors can partially cancel across features
\item
  Works well when features are "nearly" independent
\item
  Very robust to irrelevant features (they contribute near-uniform P(x\_j\textbar y=k))
\end{itemize}

\textbf{Key hyperparameters:} var\_smoothing (Gaussian NB), alpha (Multinomial/Bernoulli NB --- Laplace smoothing)

\textbf{Complexity:} Training O(n\emph{p), Inference O(p}K) where K = classes

\textbf{Strengths:} Extremely fast, works with tiny datasets, handles high-dim (text) well, good with streaming data, probabilistic output, robust to irrelevant features

\textbf{Weaknesses:} Feature independence assumption violated in most problems, poor probability calibration (use isotonic regression if you need probs), can\textquotesingle t learn feature interactions, Gaussian NB assumes normal distribution

\textbf{When to use:} Text classification (spam, sentiment), real-time classification (fast inference), small datasets, streaming/online learning

\subsubsection{4.9 k-Means Clustering}\label{49-k-means-clustering}

\paragraph{Intuition}\label{intuition-8}

Partition n data points into k clusters by assigning each point to its nearest cluster center (centroid), then updating centroids as the mean of assigned points. Iterate until convergence.

\paragraph{Lloyd\textquotesingle s Algorithm}\label{lloyds-algorithm}

\begin{verbatim}
Initialize: k centroids {mu_1, ..., mu_k} randomly (or k-means++)

Repeat until convergence:
    Assignment step:
        c_i = argmin_k ||x_i - mu_k||^2    for all i
    
    Update step:
        mu_k = (1/|C_k|) * sum_{i in C_k} x_i    for all k

Converges when assignments stop changing
\end{verbatim}

\textbf{Objective (minimized):}

\begin{verbatim}
J = sum_i ||x_i - mu_{c_i}||^2    (Within-Cluster Sum of Squares, WCSS)
\end{verbatim}

This is a non-convex optimization --- \textbf{local optima exist}. Solution: run multiple random initializations, pick the best.

\paragraph{k-Means++ Initialization}\label{k-means-initialization}

Instead of random initialization, choose centroids with probability proportional to distance from existing centroids:

\begin{verbatim}
P(x is chosen as next centroid) ∝ min_j ||x - mu_j||^2
\end{verbatim}

This spreads initial centroids far apart, typically giving better solutions and faster convergence.

\paragraph{Choosing k: The Elbow Method}\label{choosing-k-the-elbow-method}

\begin{verbatim}
Plot WCSS vs k:
        WCSS
         |
   High  |*
         | *
         |  *
         |   *___________
         |                (plateau)
    Low  +----------→ k
              ^ Elbow = good k
\end{verbatim}

The elbow is where adding more clusters has diminishing returns. Formal methods: Silhouette score, Gap statistic, BIC.

\textbf{Silhouette score:}

\begin{verbatim}
s(i) = (b(i) - a(i)) / max(a(i), b(i))
a(i) = mean intra-cluster distance
b(i) = mean nearest-cluster distance
Range: [-1, 1], higher is better
\end{verbatim}

\paragraph{Assumptions and Limitations}\label{assumptions-and-limitations}

\begin{itemize}
\tightlist
\item
  Assumes clusters are \textbf{spherical} and roughly \textbf{equal size}
\item
  Assumes Euclidean distance is meaningful
\item
  Must specify k in advance
\item
  Sensitive to outliers (use k-Medoids/PAM instead)
\item
  Can fail on ring/crescent/elongated shapes (use DBSCAN/GMM)
\end{itemize}

\textbf{Key hyperparameters:} n\_clusters (k), init (k-means++ or random), n\_init (number of restarts), max\_iter, random\_state

\textbf{Complexity:} Training O(n * k * p * iter), Inference O(k * p) per point

\subsubsection{4.10 Hierarchical Clustering}\label{410-hierarchical-clustering}

\paragraph{Intuition}\label{intuition-9}

Build a tree of clusters (dendrogram) either \textbf{bottom-up (agglomerative)} or top-down (divisive). Cut the tree at the desired level to get clusters.

\paragraph{Agglomerative Clustering}\label{agglomerative-clustering}

\begin{verbatim}
Initialize: Each point is its own cluster (n clusters)

Repeat until one cluster remains:
    Find the two closest clusters C_i and C_j
    Merge them into one cluster
    Update distance matrix

Result: Dendrogram (binary tree)
Cut at height h → determines number of clusters
\end{verbatim}

\paragraph{Linkage Criteria (how to measure cluster distance)}\label{linkage-criteria-how-to-measure-cluster-distance}

\begin{longtable}[]{@{}
  >{\raggedright\arraybackslash}p{(\columnwidth - 4\tabcolsep) * \real{0.0964}}
  >{\raggedright\arraybackslash}p{(\columnwidth - 4\tabcolsep) * \real{0.4844}}
  >{\raggedright\arraybackslash}p{(\columnwidth - 4\tabcolsep) * \real{0.4192}}@{}}
\toprule\noalign{}
\begin{minipage}[b]{\linewidth}\raggedright
Linkage
\end{minipage} & \begin{minipage}[b]{\linewidth}\raggedright
Formula
\end{minipage} & \begin{minipage}[b]{\linewidth}\raggedright
Behavior
\end{minipage} \\
\midrule\noalign{}
\endhead
\bottomrule\noalign{}
\endlastfoot
\textbf{Single} & \texttt{min\ d(a,b),\ a\ in\ C\_i,\ b\ in\ C\_j} & Chaining effect, elongated clusters \\
\textbf{Complete} & \texttt{max\ d(a,b),\ a\ in\ C\_i,\ b\ in\ C\_j} & Compact, equal-size clusters \\
\textbf{Average} & \texttt{mean\ d(a,b),\ a\ in\ C\_i,\ b\ in\ C\_j} & Balance between single and complete \\
\textbf{Ward} & Minimize within-cluster variance increase & Usually best, similar to k-means \\
\end{longtable}

\textbf{Ward linkage} is the default choice --- minimizes the increase in total WCSS when merging, producing compact, similarly-sized clusters.

\textbf{Complexity:} Naive O(n\^{}3), optimized O(n\^{}2 log n) with heap; Inference O(1) after dendrogram built

\textbf{Strengths:} No need to specify k in advance, produces dendrogram (full hierarchy), deterministic, works with any distance metric, can find non-spherical clusters (with single linkage)

\textbf{Weaknesses:} O(n\^{}2) memory for distance matrix, slow for large datasets, can\textquotesingle t easily update with new data, cutting the dendrogram is subjective

\textbf{When to use:} When you need a hierarchy, small-to-medium datasets, exploratory analysis, gene expression analysis (bioinformatics standard)

\subsubsection{4.11 DBSCAN}\label{411-dbscan}

\paragraph{Intuition}\label{intuition-10}

\textbf{Density-Based Spatial Clustering of Applications with Noise.} Clusters are regions of high point density. Points in low-density regions are labeled as \textbf{noise/outliers}. No need to specify k --- clusters emerge from density.

\paragraph{Core Concepts}\label{core-concepts}

\begin{verbatim}
epsilon (eps): neighborhood radius
min_samples: minimum points to form a dense region

Point types:
- Core point:    has >= min_samples points within eps radius (including itself)
- Border point:  within eps of a core point but < min_samples in its own eps
- Noise point:   neither core nor border (outlier)
\end{verbatim}

\paragraph{The Algorithm}\label{the-algorithm-1}

\begin{verbatim}
for each unvisited point p:
    neighbors = all points within eps of p
    
    if |neighbors| < min_samples:
        label p as NOISE
    else:
        create new cluster C
        add p to C
        expandCluster(p, neighbors, C, eps, min_samples)

function expandCluster(p, neighbors, C, eps, min_samples):
    for each q in neighbors:
        if q is NOISE: add q to C (border point)
        if q is unvisited:
            mark q as visited
            q_neighbors = all points within eps of q
            if |q_neighbors| >= min_samples:  # q is also a core point
                neighbors = neighbors U q_neighbors  # expand search
            if q not in any cluster:
                add q to C
\end{verbatim}

\textbf{Key insight:} DBSCAN can find \textbf{arbitrary-shaped clusters} --- it connects core points that are within eps of each other. It naturally handles outliers (noise points).

\paragraph{Choosing eps and min\_samples}\label{choosing-eps-and-min_samples}

\begin{itemize}
\tightlist
\item
  \textbf{min\_samples}: Rule of thumb: \texttt{min\_samples\ \textgreater{}=\ dimensionality\ +\ 1}, usually 4-10. Higher = less noise sensitivity.
\item
  \textbf{eps}: Plot k-distance graph (distance to k-th nearest neighbor for each point), look for the "elbow." The knee of the curve is a good eps.
\end{itemize}

\textbf{Complexity:} O(n log n) with spatial indexing (kd-tree/ball-tree), O(n\^{}2) naively

\textbf{Strengths:} Finds arbitrary shapes, automatically detects outliers, no need to specify k, robust to different cluster sizes

\textbf{Weaknesses:} Struggles with varying densities (HDBSCAN solves this), sensitive to eps and min\_samples, doesn\textquotesingle t work well in high dimensions (epsilon loses meaning), not good at classifying new points

\textbf{When to use:} Geospatial clustering, anomaly/outlier detection, when you expect non-spherical clusters, when number of clusters is unknown

\subsubsection{4.12 PCA / Dimensionality Reduction}\label{412-pca--dimensionality-reduction}

\paragraph{Intuition}\label{intuition-11}

Find the \textbf{directions of maximum variance} in the data and project onto fewer dimensions. The first principal component explains the most variance, the second (orthogonal to first) explains the second most, etc.

\paragraph{The Math}\label{the-math-2}

\textbf{Setup:} X is n x p, centered (subtract mean from each feature)

\textbf{Covariance matrix:}

\begin{verbatim}
Sigma = (1/n) * X^T X     (p x p matrix)
\end{verbatim}

\textbf{Eigenvector decomposition:}

\begin{verbatim}
Sigma = V Lambda V^T

V = [v_1 | v_2 | ... | v_p]    (eigenvectors = principal component directions)
Lambda = diag(lambda_1, lambda_2, ..., lambda_p)   (eigenvalues, sorted descending)
\end{verbatim}

\textbf{Project onto k components:}

\begin{verbatim}
Z = X * V_k    (n x k matrix)
V_k = first k columns of V
\end{verbatim}

\textbf{Explained variance ratio:}

\begin{verbatim}
EVR_j = lambda_j / sum_i lambda_i
Cumulative EVR: choose k such that sum_{j=1}^k EVR_j >= 0.95 (95% threshold)
\end{verbatim}

\paragraph{SVD Connection}\label{svd-connection}

\begin{verbatim}
X = U Sigma V^T    (SVD)
Principal components = right singular vectors V
Scores = U Sigma = X V
Eigenvalues of covariance = sigma_i^2 / n (squared singular values)
\end{verbatim}

In practice, \textbf{always use SVD, not eigendecomposition} --- more numerically stable, works even when p \textgreater\textgreater{} n.

\paragraph{Geometric Intuition}\label{geometric-intuition}

\begin{verbatim}
Original 2D data:          After PCA:
     *  * *                PC1: direction of max spread
   *  * * *              /
  * * * *           ----/-------→ PC1 (new x-axis)
   * * *               /
     *              PC2: orthogonal, less spread
                       ↑ (new y-axis)
\end{verbatim}

The first PC is the axis along which projecting the data retains the most information (minimum reconstruction error).

\paragraph{PCA Assumptions}\label{pca-assumptions}

\begin{itemize}
\tightlist
\item
  Directions of max \textbf{variance} = directions of interest (this fails when important patterns are in low-variance directions)
\item
  Linear relationships between features
\item
  Data is centered (subtract mean --- PCA is not scale-invariant, must standardize)
\end{itemize}

\paragraph{Kernel PCA}\label{kernel-pca}

Apply PCA in the kernel-transformed feature space:

\begin{verbatim}
Replace X^T X with kernel matrix K where K_{ij} = k(x_i, x_j)
\end{verbatim}

Finds non-linear low-dimensional structure.

\paragraph{t-SNE (brief)}\label{t-sne-brief}

For \textbf{visualization} only (2D/3D). Preserves local structure --- nearby points in high-dim stay nearby in low-dim. Non-parametric, can\textquotesingle t transform new points. Use PCA for preprocessing before t-SNE on large datasets.

\paragraph{UMAP (brief)}\label{umap-brief}

Faster than t-SNE, preserves more global structure. Can transform new points. Preferred for large-scale visualization and as a preprocessing step.

\textbf{Key hyperparameters (PCA):} n\_components, svd\_solver (auto, full, randomized, arpack), whiten

\textbf{Complexity:} O(n * p\^{}2 + p\^{}3) exact, O(n * p * k) randomized SVD

\textbf{When to use PCA:}

\begin{itemize}
\tightlist
\item
  Remove multicollinearity before linear models
\item
  Speed up training (reduce features)
\item
  Visualization (first 2-3 components)
\item
  Noise reduction
\item
  Feature engineering for downstream models
\end{itemize}

\subsubsection{4.13 Ensemble Methods: Bagging vs Boosting vs Stacking}\label{413-ensemble-methods-bagging-vs-boosting-vs-stacking}

\paragraph{Bagging (Bootstrap AGGregating)}\label{bagging-bootstrap-aggregating}

\begin{verbatim}
Core idea: Reduce variance by averaging many high-variance models
Training: Independent, parallel -- models don't see each other's errors
Sampling: Bootstrap (with replacement)
Combination: Average (regression), majority vote (classification)
\end{verbatim}

\begin{verbatim}
Data
 |
 +----Bootstrap_1 ---→ Model_1 ----+
 |                                  |
 +----Bootstrap_2 ---→ Model_2 ----+--→ Average/Vote ---→ Final Prediction
 |                                  |
 +----Bootstrap_3 ---→ Model_3 ----+
\end{verbatim}

\textbf{Examples:} Random Forest (bagging + feature randomness), Bagged SVMs

\textbf{Best for:} High-variance, low-bias base models (deep trees)

\paragraph{Boosting}\label{boosting}

\begin{verbatim}
Core idea: Reduce bias by sequentially correcting errors
Training: Sequential -- each model focuses on previous model's mistakes
Sampling: Full dataset (or weighted/modified)
Combination: Weighted sum (adaptive weighting)
\end{verbatim}

\begin{verbatim}
Data → Model_1 → Errors
                   |
              Model_2 (focus on errors of Model_1) → Residuals
                                                          |
                                                     Model_3 (focus on Model_2's residuals)
                                                                ...
                                                          Final = weighted sum
\end{verbatim}

\textbf{Examples:} AdaBoost, GBM, XGBoost, LightGBM

\textbf{Best for:} Low-variance, high-bias base models (shallow trees, stumps)

\paragraph{Stacking (Stacked Generalization)}\label{stacking-stacked-generalization}

\begin{verbatim}
Core idea: Use a meta-model to optimally combine base model predictions
Training: Two-level training with cross-validation to avoid leakage
\end{verbatim}

\begin{verbatim}
Level 0 (Base models):
  X ---→ Model_A (Logistic Regression)  --→ pred_A ----+
  X ---→ Model_B (Random Forest)        --→ pred_B ----+--→ [pred_A, pred_B, pred_C]
  X ---→ Model_C (XGBoost)              --→ pred_C ----+

Level 1 (Meta-model):
  [pred_A, pred_B, pred_C] ---→ Meta Learner (often Logistic Regression) ---→ Final
\end{verbatim}

\textbf{Stacking implementation (preventing leakage):}

\begin{Shaded}
\begin{Highlighting}[]
\CommentTok{\# Cross{-}val approach}
\NormalTok{oof\_predictions }\OperatorTok{=}\NormalTok{ np.zeros((n\_train, n\_models))}
\ControlFlowTok{for}\NormalTok{ model\_idx, model }\KeywordTok{in} \BuiltInTok{enumerate}\NormalTok{(base\_models):}
    \ControlFlowTok{for}\NormalTok{ fold\_idx, (train\_idx, val\_idx) }\KeywordTok{in} \BuiltInTok{enumerate}\NormalTok{(kfold.split(X)):}
\NormalTok{        model.fit(X[train\_idx], y[train\_idx])}
\NormalTok{        oof\_predictions[val\_idx, model\_idx] }\OperatorTok{=}\NormalTok{ model.predict(X[val\_idx])}

\CommentTok{\# Train meta{-}model on out{-}of{-}fold predictions}
\NormalTok{meta\_model.fit(oof\_predictions, y)}
\end{Highlighting}
\end{Shaded}

\textbf{Strengths:} Often best accuracy; Weaknesses: complex, slow, overfitting risk if leakage

\paragraph{Comparison Table}\label{comparison-table-1}

\begin{longtable}[]{@{}
  >{\raggedright\arraybackslash}p{(\columnwidth - 6\tabcolsep) * \real{0.2135}}
  >{\raggedright\arraybackslash}p{(\columnwidth - 6\tabcolsep) * \real{0.2921}}
  >{\raggedright\arraybackslash}p{(\columnwidth - 6\tabcolsep) * \real{0.2247}}
  >{\raggedright\arraybackslash}p{(\columnwidth - 6\tabcolsep) * \real{0.2697}}@{}}
\toprule\noalign{}
\begin{minipage}[b]{\linewidth}\raggedright
Property
\end{minipage} & \begin{minipage}[b]{\linewidth}\raggedright
Bagging
\end{minipage} & \begin{minipage}[b]{\linewidth}\raggedright
Boosting
\end{minipage} & \begin{minipage}[b]{\linewidth}\raggedright
Stacking
\end{minipage} \\
\midrule\noalign{}
\endhead
\bottomrule\noalign{}
\endlastfoot
\textbf{Error type targeted} & Variance & Bias & Both \\
\textbf{Training order} & Parallel & Sequential & Two-stage \\
\textbf{Data sampling} & Bootstrap & Weighted/full & Full (CV for meta) \\
\textbf{Base model} & High-variance (deep trees) & High-bias (stumps) & Diverse \\
\textbf{Combination} & Average/vote & Weighted sum & Meta-model \\
\textbf{Speed} & Fast (parallel) & Slow (sequential) & Medium \\
\textbf{Overfitting risk} & Low & Medium (can overfit) & High (if leakage) \\
\textbf{Example} & Random Forest & XGBoost & Kaggle winning solutions \\
\textbf{When to use} & Reduce overfitting & Boost a weak model & Maximum accuracy \\
\end{longtable}

\subsection{5. Architecture / Diagrams}\label{5-architecture--diagrams}

\subsubsection{Decision Tree Structure}\label{decision-tree-structure}

\begin{verbatim}
                    [Age < 30?]                    ← Root node
                   /           \
              YES /             \ NO
                /               \
        [Income > 50k?]      [Own Home?]           ← Internal nodes
          /       \            /      \
        YES       NO         YES       NO
        /          \         /          \
   [Approve]    [Deny]  [Approve]    [Credit?]     ← Leaves + node
                                       /    \
                                     YES     NO
                                    [Approve] [Deny]

Split criterion: minimize Gini/Entropy
Pruning: remove nodes that don't significantly reduce impurity
\end{verbatim}

\subsubsection{SVM Maximum Margin Boundary}\label{svm-maximum-margin-boundary}

\begin{verbatim}
Feature 2 (x2)
     |
   6 |    ○  ○
     |       ○         margin = 2/||w||
   5 | ○    ○    - - - - - - - - -  ← w^T x + b = +1
     |          /     /
   4 |         /  ●  /             ← Decision hyperplane w^T x + b = 0
     |        / ●   /
   3 |       / ● ● /
     |      - - - - - - - - - - -  ← w^T x + b = -1
   2 |   ●
     |
   1 |__________________________ Feature 1 (x1)

 ○ = Class +1    ● = Class -1
 Support Vectors: circled points that lie ON the margin planes
 Only support vectors matter for the decision — all others can be removed
\end{verbatim}

\subsubsection{k-Means Cluster Formation}\label{k-means-cluster-formation}

\begin{verbatim}
Iteration 1: Random init        Iteration 3: Converging
   * *  ×                          * *  ×
  * * *   ×  ×         →          * *    ×  ×
 *  * *     × ×                  * *      × ×
   ■ ■  ■                           ■ ■  ■
    ■ ■                              ■ ■
 *=cluster1  ×=cluster2  ■=cluster3

Centroids (marked with +):
   + * *  ×  ×                      * + *  × ×
  * * *   ×  ×         →          * * *   ×+ ×
 *  *       × ×                  * * *     × ×
   ■ ■  ■                           ■ ■  +
    ■ ■                              ■ ■

After convergence: centroid = exact mean of each cluster
\end{verbatim}

\subsubsection{Bagging vs Boosting Flow}\label{bagging-vs-boosting-flow}

\begin{verbatim}
BAGGING (Random Forest):
┌─────────────────────────────────────────────────────────┐
│  Training Data                                           │
│       │                                                  │
│  ┌────┴────┐  ┌─────────┐  ┌─────────┐                 │
│  │Bootstrap│  │Bootstrap│  │Bootstrap│  ...             │
│  │  1      │  │   2     │  │   3     │                  │
│  └────┬────┘  └────┬────┘  └────┬────┘                 │
│       │            │            │                        │
│   ┌───▼───┐    ┌───▼───┐    ┌───▼───┐                  │
│   │ Tree1 │    │ Tree2 │    │ Tree3 │  (parallel)       │
│   └───┬───┘    └───┬───┘    └───┬───┘                  │
│       └────────────┼────────────┘                        │
│                    │                                     │
│              ┌─────▼─────┐                              │
│              │   AVERAGE  │                             │
│              │ (or vote)  │                             │
│              └─────┬─────┘                              │
│                    ▼                                     │
│              Final Prediction                           │
└─────────────────────────────────────────────────────────┘

BOOSTING (XGBoost):
┌─────────────────────────────────────────────────────────┐
│  Training Data ──────────────────────────────────────┐  │
│       │                                              │  │
│   ┌───▼───┐                                          │  │
│   │ Tree1 │                                          │  │
│   └───┬───┘                                          │  │
│       │ Residuals/Pseudo-residuals                   │  │
│   ┌───▼───┐                                          │  │
│   │ Tree2 │ (fits residuals of Tree1)                │  │
│   └───┬───┘                                          │  │
│       │ Residuals                                    │  │
│   ┌───▼───┐                                          │  │
│   │ Tree3 │ (fits residuals of Tree1+2)              │  │
│   └───┬───┘                                          │  │
│       │        SEQUENTIAL, each tree corrects        │  │
│  ┌────▼──────────────────────────────────────────┐  │  │
│  │ Final = lr*Tree1 + lr*Tree2 + lr*Tree3 + ...  │  │  │
│  └────────────────────────────────────────────────┘  │  │
└─────────────────────────────────────────────────────────┘
\end{verbatim}

\subsubsection{PCA Projection}\label{pca-projection}

\begin{verbatim}
High-Dimensional Space (2D shown):        Projected (1D PC1):
         x2
          |  ×  ×
        5 | × × ×    PC1 direction          ×  × × ×  ×
          | ×  × ×   (max variance)  →  ───×──×─×──×──×──→ PC1
        3 |   × × ×  ↗
          |  × × ×
        1 |                              Information preserved: 87% variance
          +───────→ x1                  Information lost: 13% variance

Original: 2 coordinates per point
Reduced:  1 coordinate per point (87% info kept)

Eigenvectors: orthogonal axes of new coordinate system
Eigenvalues: amount of variance along each axis
\end{verbatim}

\subsubsection{Bias-Variance: Single Tree vs Random Forest}\label{bias-variance-single-tree-vs-random-forest}

\begin{verbatim}
Error
  |
  |  Single Tree:
  |  ●─────────────────●  Total Error
  |  |   high variance  |
  |  └──── bias (low) ──┘
  |
  |  Random Forest:
  |  ●──────────────────●  Total Error (lower!)
  |  |  lower variance   |
  |  └─── bias (same) ───┘
  |
  +─────────────────────────→ Model Complexity

Variance: |────────| (big for single tree) → |────| (small for RF)
Bias:     |──| (small for both deep trees)
Noise:    constant

RF total error < Single Tree total error
\end{verbatim}

\subsection{6. Real-World Examples}\label{6-real-world-examples}

\subsubsection{Linear Regression}\label{linear-regression}

\begin{itemize}
\tightlist
\item
  \textbf{House price prediction} (Zillow Zestimate started as linear regression with 200+ features)
\item
  \textbf{Stock return forecasting} (Fama-French factor models)
\item
  \textbf{A/B test analysis} (linear models for CUPED variance reduction at Netflix)
\end{itemize}

\subsubsection{Logistic Regression}\label{logistic-regression}

\begin{itemize}
\tightlist
\item
  \textbf{Credit scoring} (FICO score model, FDIC regulatory requirement for interpretability)
\item
  \textbf{Email spam} (Google\textquotesingle s early Spam filter)
\item
  \textbf{Google Ads CTR} (billion-scale with FTRL optimization)
\end{itemize}

\subsubsection{Decision Trees}\label{decision-trees}

\begin{itemize}
\tightlist
\item
  \textbf{Medical diagnosis protocols} (sepsis screening, CHA2DS2-VASc stroke risk)
\item
  \textbf{Customer churn rules} (easily communicated to business stakeholders)
\item
  \textbf{Loan approval} (regulatory requirement for explainability in EU AI Act)
\end{itemize}

\subsubsection{Random Forest}\label{random-forest}

\begin{itemize}
\tightlist
\item
  \textbf{Genomics} (predicting gene expression from SNPs --- thousands of features)
\item
  \textbf{Fraud detection at PayPal} (100+ features, handles missing data naturally)
\item
  \textbf{Remote sensing} (classifying satellite pixels into land-use categories)
\end{itemize}

\subsubsection{Gradient Boosting (XGBoost/LightGBM)}\label{gradient-boosting-xgboostlightgbm}

\begin{itemize}
\tightlist
\item
  \textbf{Kaggle}: Won \textgreater70\% of structured data competitions since 2016
\item
  \textbf{LinkedIn Feed Ranking}: LightGBM for post relevance scoring
\item
  \textbf{Uber surge pricing}: GBM for demand forecasting
\item
  \textbf{Microsoft Bing}: LightGBM for document relevance ranking
\end{itemize}

\subsubsection{SVM}\label{svm}

\begin{itemize}
\tightlist
\item
  \textbf{OCR (handwriting recognition)} --- original MNIST benchmark
\item
  \textbf{Bioinformatics}: Protein classification, gene expression microarrays (high-dim, small n)
\item
  \textbf{Face detection}: Early face detectors used SVMs with Haar features
\end{itemize}

\subsubsection{k-NN}\label{k-nn}

\begin{itemize}
\tightlist
\item
  \textbf{Netflix/Spotify recommendations} --- k-NN in embedding space (collaborative filtering)
\item
  \textbf{Medical diagnosis from embeddings}: "Patients similar to this one had X outcome"
\item
  \textbf{Fraud detection}: k-NN anomaly scores
\end{itemize}

\subsubsection{Naive Bayes}\label{naive-bayes}

\begin{itemize}
\tightlist
\item
  \textbf{Spam filters}: SpamAssassin, iCloud Mail
\item
  \textbf{Sentiment analysis}: Fast, surprisingly effective baseline
\item
  \textbf{Medical text classification}: On-device, privacy-preserving
\end{itemize}

\subsubsection{k-Means}\label{k-means}

\begin{itemize}
\tightlist
\item
  \textbf{Customer segmentation}: RFM (Recency, Frequency, Monetary) clusters for marketing
\item
  \textbf{Image compression}: Color quantization (k-Means on RGB values \(\rightarrow\) k representative colors)
\item
  \textbf{Document clustering}: TF-IDF vectors \(\rightarrow\) topic groups
\end{itemize}

\subsubsection{PCA}\label{pca}

\begin{itemize}
\tightlist
\item
  \textbf{Eigenfaces}: Early face recognition (Turk \& Pentland, 1991)
\item
  \textbf{Risk factor models in finance}: 5-10 PCs explain 70\%+ of stock return variance
\item
  \textbf{NLP preprocessing}: PCA on TF-IDF matrix before SVM
\end{itemize}

\subsection{7. Real-Life Analogies}\label{7-real-life-analogies}

\emph{Imagine a hospital emergency room. A \textbf{panel of specialist doctors} uses different methods to diagnose patients --- each algorithm is a doctor with a distinct diagnostic style.}

\begin{longtable}[]{@{}
  >{\raggedright\arraybackslash}p{(\columnwidth - 4\tabcolsep) * \real{0.1211}}
  >{\raggedright\arraybackslash}p{(\columnwidth - 4\tabcolsep) * \real{0.5354}}
  >{\raggedright\arraybackslash}p{(\columnwidth - 4\tabcolsep) * \real{0.3434}}@{}}
\toprule\noalign{}
\begin{minipage}[b]{\linewidth}\raggedright
Algorithm
\end{minipage} & \begin{minipage}[b]{\linewidth}\raggedright
Doctor Analogy
\end{minipage} & \begin{minipage}[b]{\linewidth}\raggedright
Key Behavior
\end{minipage} \\
\midrule\noalign{}
\endhead
\bottomrule\noalign{}
\endlastfoot
\textbf{Decision Tree} & The triage nurse with a printed flowchart: "Temp \textgreater{} 38.5°C? \(\rightarrow\) Yes \(\rightarrow\) Chest pain? \(\rightarrow\) No \(\rightarrow\) Rule out pneumonia." & Follows explicit branching rules; fast but rigid; one wrong branch goes down the wrong path \\
\textbf{Random Forest} & The \textbf{medical board}: 100 doctors each examine a random patient subset and a random subset of symptoms; majority vote wins & Cancels out individual doctor errors; much more reliable than any single doctor \\
\textbf{Gradient Boosting} & The \textbf{residency training program}: Intern makes a diagnosis \(\rightarrow\) Senior corrects the mistakes \(\rightarrow\) Fellow corrects remaining errors \(\rightarrow\) Sequential improvement & Each successive doctor specializes in what the previous missed; powerful but takes time to train \\
\textbf{SVM} & The \textbf{surgeon} who draws the clearest anatomical line between healthy and diseased tissue, maximizing the margin for error & Maximizes separation; only cares about boundary cases (support vectors) \\
\textbf{k-NN} & The doctor who says \emph{"Let me look up the 5 most similar patients we\textquotesingle ve seen and see what happened to them"} & No model training --- pure case-based reasoning; only as good as the patient database \\
\textbf{Naive Bayes} & The \textbf{phone triage doctor} who assumes symptoms are independent: "Fever? +10 pts. Cough? +5 pts. Fatigue? +3 pts." & Fast; wrong about independence (fever and cough often co-occur) but often right enough \\
\textbf{Logistic Regression} & The \textbf{general practitioner} who weighs each symptom linearly: "I assign 2.3\(\times\) weight to chest pain and 0.7\(\times\) to age" and gives a probability & Interpretable weights; can explain decision; good calibrated probabilities \\
\textbf{k-Means} & The \textbf{epidemiologist} grouping patients into symptom clusters: "Group A = respiratory, Group B = cardiovascular" & Discovers natural groupings without labels; must specify number of groups upfront \\
\textbf{PCA} & The \textbf{chief diagnostician} who identifies the 3 "vital sign combinations" (principal components) that explain 95\% of patient variation & Reduces 50 measurements to 3 key dimensions; loses some detail but gains clarity \\
\textbf{DBSCAN} & The \textbf{outbreak investigator} finding clusters of patients who live near each other + flagging isolated cases as "outliers" & No need to specify clusters; naturally identifies anomalies \\
\textbf{Hierarchical Clustering} & The \textbf{taxonomy specialist} building a family tree of related conditions from common to rare & Produces a full hierarchy; shows how conditions relate at different levels of similarity \\
\textbf{Stacking} & The \textbf{hospital committee}: GP, cardiologist, and pulmonologist each give a diagnosis; a chief of medicine combines their opinions with learned weights & Best accuracy by leveraging diverse specialist views; but complex and slow \\
\end{longtable}

\emph{The best hospitals (and ML systems) use the right specialist for the right case --- and often consult the full board for the hardest diagnoses.}

\subsection{8. Memory Tricks / Mnemonics}\label{8-memory-tricks--mnemonics}

\subsubsection{Algorithm Selection: "GLDRSNK-BPS" (Great Lakes Don\textquotesingle t Run South, Never Kite --- Big Pandas Swim)}\label{algorithm-selection-gldrsnk-bps-great-lakes-dont-run-south-never-kite--big-pandas-swim}

\begin{itemize}
\tightlist
\item
  \textbf{G}radient Boosting: tabular, winning solution
\item
  \textbf{L}ogistic Regression: baseline, interpretable, fast
\item
  \textbf{D}ecision Tree: interpretable, rules
\item
  \textbf{R}andom Forest: robust, high-dim, OOB
\item
  \textbf{S}VM: small n, large p, kernels
\item
  \textbf{N}aive Bayes: text, speed, streaming
\item
  \textbf{K}-NN: recommendation, embeddings
\item
  \textbf{B}oosting (GBM/XGB): sequential correction
\item
  \textbf{P}CA: dimensionality, visualization
\item
  \textbf{S}tacking: Kaggle, maximum accuracy
\end{itemize}

\subsubsection{Key Formulas --- SSGH (Signed Symmetrical, Gradient, Hinge)}\label{key-formulas--ssgh-signed-symmetrical-gradient-hinge}

\begin{itemize}
\tightlist
\item
  \textbf{Sigmoid:} "S-shape squeezes infinity into {[}0,1{]}"
\item
  \textbf{Softmax:} "Sigmoid\textquotesingle s multiclass cousin --- normalizes exponentials"
\item
  \textbf{Gini:} "1 minus sum of squares" \(\rightarrow\) \texttt{1\ -\ sum(p\^{}2)}
\item
  \textbf{Hinge:} "max(0, 1-y*f(x))" \(\rightarrow\) zero for correct confident predictions
\end{itemize}

\subsubsection{Bias-Variance Mnemonic: "OVER-UNDER-JUST"}\label{bias-variance-mnemonic-over-under-just}

\begin{itemize}
\tightlist
\item
  \textbf{Overfit} = High Variance, Low Bias (deep tree, k-NN with k=1)
\item
  \textbf{Underfit} = High Bias, Low Variance (linear model on curved data)
\item
  \textbf{Just right} = Bagging reduces variance; Boosting reduces bias
\end{itemize}

\subsubsection{Kernel Trick: "Measure similarity, never transform"}\label{kernel-trick-measure-similarity-never-transform}

\begin{itemize}
\tightlist
\item
  "You don\textquotesingle t need to go to infinite dimensions --- just measure how similar points are in that space"
\item
  \texttt{K(x,z)} = inner product in feature space without computing feature vectors
\end{itemize}

\subsubsection{Regularization Directions:}\label{regularization-directions}

\begin{itemize}
\tightlist
\item
  \textbf{Ridge (L2):} "Ridge = Round ball \(\rightarrow\) weights shrink uniformly toward zero, never exactly zero"
\item
  \textbf{Lasso (L1):} "Lasso = Diamond/Spiky \(\rightarrow\) corners force weights to exactly zero \(\rightarrow\) sparse \(\rightarrow\) selection"
\end{itemize}

\subsubsection{k-Means vs DBSCAN:}\label{k-means-vs-dbscan}

\begin{itemize}
\tightlist
\item
  "k-Means needs k, DBSCAN finds k"
\item
  "k-Means finds spheres, DBSCAN finds shapes"
\item
  "k-Means hates outliers, DBSCAN loves to find them"
\end{itemize}

\subsubsection{PCA: "Eigenvectors point where data varies most, eigenvalues say how much"}\label{pca-eigenvectors-point-where-data-varies-most-eigenvalues-say-how-much}

\subsubsection{\texorpdfstring{Ensemble Mnemonic: "Bag \(\rightarrow\) Parallel, Boost \(\rightarrow\) Sequential, Stack \(\rightarrow\) Hierarchical"}{Ensemble Mnemonic: "Bag \textbackslash rightarrow Parallel, Boost \textbackslash rightarrow Sequential, Stack \textbackslash rightarrow Hierarchical"}}\label{ensemble-mnemonic-bag-rightarrow-parallel-boost-rightarrow-sequential-stack-rightarrow-hierarchical}

\subsection{9. Common Interview Questions}\label{9-common-interview-questions}

\subsubsection{Linear Regression}\label{linear-regression-1}

\textbf{Q: What are the assumptions of linear regression and what happens when they\textquotesingle re violated?}

\textbf{Model Answer:}
"There are 5 key assumptions (Gauss-Markov conditions):

\begin{enumerate}
\def\labelenumi{\arabic{enumi}.}
\tightlist
\item
  \textbf{Linearity}: \texttt{E{[}y\textbar{}X{]}\ =\ Xbeta}. Violation \(\rightarrow\) biased predictions. Fix: add polynomial features, use trees.
\item
  \textbf{No multicollinearity}: features not perfectly correlated. Violation \(\rightarrow\) unstable coefficients, wrong signs. Fix: Ridge regression, PCA, drop correlated features.
\item
  \textbf{Homoscedasticity}: constant error variance. Violation \(\rightarrow\) inefficient OLS estimates, wrong standard errors. Fix: log-transform y, weighted least squares, robust SEs.
\item
  \textbf{No autocorrelation}: residuals uncorrelated. Violation in time-series \(\rightarrow\) underestimated variance. Fix: ARIMA, GLS.
\item
  \textbf{Exogeneity}: errors uncorrelated with features. Violation \(\rightarrow\) biased estimates. Fix: instrumental variables.
  For inference we also need normality of residuals, but OLS predictions remain unbiased without it."
\end{enumerate}

\textbf{Follow-up:} "How would you detect heteroscedasticity?"
"Plot residuals vs. fitted values --- if variance increases with fitted values, it\textquotesingle s heteroscedastic. Breusch-Pagan test formally."

\textbf{Q: When would you use Ridge vs Lasso?}

\textbf{Model Answer:}
"Ridge (L2) when you believe all features are somewhat relevant --- it shrinks all coefficients proportionally without zeroing any out. Good for multicollinearity --- the added lambda*I term to \texttt{X\^{}T\ X} ensures invertibility.
Lasso (L1) when you want automatic feature selection --- it produces sparse solutions (exact zeros) because the L1 ball has corners at the axes; the constraint boundary tends to intersect at corners. Better when true model is sparse.
ElasticNet when you have correlated features AND want sparsity --- Lasso tends to arbitrarily pick one feature from a correlated group; ElasticNet selects the group."

\subsubsection{Decision Trees \& Ensembles}\label{decision-trees--ensembles}

\textbf{Q: Why does Random Forest work? Why is it better than a single decision tree?}

\textbf{Model Answer:}
"Decision trees have \textbf{high variance} --- small changes in training data produce very different trees. Random Forest exploits the bias-variance decomposition. If you average \texttt{B} trees with variance \texttt{sigma\^{}2} and pairwise correlation \texttt{rho}, the ensemble variance is \texttt{rho*sigma\^{}2\ +\ (1-rho)/B\ *\ sigma\^{}2}. By bootstrap sampling and random feature selection, we reduce \texttt{rho} (decorrelate the trees) without increasing bias. As \texttt{B} grows, the first term dominates --- \texttt{rho*sigma\^{}2} --- which is why more trees always helps up to a point, but returns diminish once you\textquotesingle ve covered the correlation structure."

\textbf{Follow-up:} "What is OOB error? How is it different from cross-validation?"
"Each bootstrap sample excludes \textasciitilde36.8\% of training data. These out-of-bag samples serve as a natural validation set for each tree. Average OOB error across trees estimates test error without needing a separate validation set. It\textquotesingle s essentially built-in \textquotesingle leave-\textasciitilde37\%-out\textquotesingle{} cross-validation. The difference from k-fold CV: OOB uses the model trained on the bootstrap sample (not fold), so each tree is evaluated on \textasciitilde37\% of data rather than exactly 1/k."

\textbf{Q: Explain gradient boosting from first principles.}

\textbf{Model Answer:}
"Gradient boosting is \textbf{gradient descent in function space}. We want to minimize some loss \texttt{L(y,\ F(x))}. Instead of optimizing weights \texttt{w} (as in neural networks), we optimize the function \texttt{F} itself.

Step 1: Start with a constant prediction: \texttt{F\_0(x)\ =\ argmin\ sum\ L(y\_i,\ c)}.
Step 2: Compute pseudo-residuals = negative gradient of loss w.r.t. current predictions:
\texttt{r\_\{im\}\ =\ -{[}dL(y\_i,\ F(x\_i))/dF(x\_i){]}}
For MSE, this is just \texttt{y\_i\ -\ F(x\_i)} --- the regular residuals.
Step 3: Fit a shallow tree to these pseudo-residuals.
Step 4: Add this tree to the ensemble with a step size (learning rate).
Step 5: Repeat.

The insight is that each tree is a \textbf{weak correction step} in function space. The sequence of trees collectively performs gradient descent on the loss function. XGBoost improves this by adding tree complexity regularization directly to the objective, approximating the optimal tree structure via a second-order Taylor expansion of the loss."

\textbf{Follow-up:} "What\textquotesingle s the difference between XGBoost and LightGBM?"
"Two key differences: (1) Tree growth strategy --- XGBoost grows level-wise (all nodes at a depth before going deeper), LightGBM grows leaf-wise (always splits the leaf with max delta loss). Leaf-wise is more accurate with fewer trees but needs \texttt{num\_leaves} control to avoid overfitting. (2) Data handling --- LightGBM uses histogram binning for features (reduces memory and speeds up splits), GOSS sampling (keep high-gradient points, subsample low-gradient), and EFB feature bundling. Combined, LightGBM is 10-100x faster than XGBoost on large datasets."

\subsubsection{SVM}\label{svm-1}

\textbf{Q: Explain the kernel trick. Why can\textquotesingle t we just add polynomial features manually?}

\textbf{Model Answer:}
"The kernel trick lets us implicitly map data to a higher (or infinite) dimensional space without ever computing the feature vectors in that space.

Manually adding polynomial features: for degree-d polynomial on p features, you get \texttt{O(p\^{}d)} features --- for p=100, d=3, that\textquotesingle s millions. This is computationally infeasible and causes memory issues.

The kernel trick works because the SVM dual formulation only depends on data through \textbf{inner products} \texttt{x\_i\^{}T\ x\_j}. If we have a function \texttt{K(x\_i,\ x\_j)\ =\ phi(x\_i)\^{}T\ phi(x\_j)} that computes the inner product in the high-dim space \textbf{directly}, we never need to compute \texttt{phi(x)} explicitly.

For the RBF/Gaussian kernel: \texttt{K(x,z)\ =\ exp(-gamma\textbar{}\textbar{}x-z\textbar{}\textbar{}\^{}2)}, the implicit feature space is \textbf{infinite-dimensional}. We\textquotesingle re essentially projecting to a space where any reasonable data can be linearly separated. But we never pay the computational cost of that infinite-dimensional space --- kernel evaluation is O(p), not O(\(\infty\))."

\textbf{Follow-up:} "When would you NOT use an SVM?"
"When n \textgreater{} \textasciitilde50,000 --- training is O(n\^{}2) to O(n\^{}3). When you need probability outputs (SVMs don\textquotesingle t give them; Platt scaling is a post-hoc hack). When features are highly correlated and numerous. In these cases, gradient boosting or neural networks are better choices."

\subsubsection{Clustering}\label{clustering}

\textbf{Q: How do you choose k in k-Means? How do you evaluate clustering quality?}

\textbf{Model Answer:}
"For choosing k:

\begin{itemize}
\tightlist
\item
  \textbf{Elbow method}: Plot WCSS vs k; the \textquotesingle elbow\textquotesingle{} is where adding more clusters has diminishing returns. Subjective and often ambiguous.
\item
  \textbf{Silhouette score}: For each point, compute \texttt{(b\ -\ a)/max(a,b)} where \texttt{a} = mean intra-cluster distance, \texttt{b} = mean nearest-cluster distance. Range {[}-1,1{]}, higher is better. Average across all points and choose k that maximizes it.
\item
  \textbf{Gap statistic}: Compare WCSS to expected WCSS under null distribution (uniform random data). Choose k where the gap is largest.
\item
  \textbf{Domain knowledge}: Often the most practical --- business requirements define the number of segments.
\end{itemize}

For evaluation:

\begin{itemize}
\tightlist
\item
  \textbf{Internal metrics} (no ground truth): Silhouette, Davies-Bouldin index, Calinski-Harabasz
\item
  \textbf{External metrics} (when ground truth exists): Adjusted Rand Index (ARI), Normalized Mutual Information (NMI), homogeneity, completeness"
\end{itemize}

\subsubsection{PCA}\label{pca-1}

\textbf{Q: Explain PCA from first principles. Why is it useful?}

\textbf{Model Answer:}
"PCA finds directions of maximum variance in the data. Mathematically:

\begin{enumerate}
\def\labelenumi{\arabic{enumi}.}
\tightlist
\item
  Center the data: \texttt{X\_c\ =\ X\ -\ mean(X)}
\item
  Compute covariance matrix: \texttt{Sigma\ =\ (1/n)\ X\_c\^{}T\ X\_c}
\item
  Eigen-decompose: \texttt{Sigma\ =\ V\ Lambda\ V\^{}T} where V\textquotesingle s columns are principal components
\item
  Project: \texttt{Z\ =\ X\_c\ V\_k} using the top k eigenvectors
\end{enumerate}

The first PC is the direction that explains the most variance in the data. Second PC is orthogonal to first and explains the second most variance. By projecting onto top-k PCs, we retain most information in fewer dimensions.

Why useful:

\begin{itemize}
\tightlist
\item
  \textbf{Removes multicollinearity} --- PCs are orthogonal by construction
\item
  \textbf{Speeds up training} --- fewer features
\item
  \textbf{Denoises} --- noise often lives in low-variance components
\item
  \textbf{Visualization} --- project to 2D/3D
\end{itemize}

Important caveats: PCA assumes variance = importance. If important features have low variance, PCA will discard them. Also, PCs are linear combinations of features --- interpretability is lost. And PCA is scale-sensitive --- always standardize features first."

\subsubsection{Bias-Variance}\label{bias-variance}

\textbf{Q: A model performs well on training data but poorly on test data. What\textquotesingle s wrong, and how do you fix it?}

\textbf{Model Answer:}
"This is overfitting --- high variance, low bias. The model has memorized training data including noise.

Diagnosis: large gap between train error and validation error.

Fixes in order of what to try:

\begin{enumerate}
\def\labelenumi{\arabic{enumi}.}
\tightlist
\item
  \textbf{More data} (best fix --- reduces variance directly)
\item
  \textbf{Regularization} (L1/L2 for linear models, max\_depth/min\_samples for trees, dropout for neural nets)
\item
  \textbf{Simpler model} (less expressive model family)
\item
  \textbf{Feature selection} (reduce dimensionality)
\item
  \textbf{Ensemble methods} (bagging reduces variance)
\item
  \textbf{Cross-validation} (ensure train/test aren\textquotesingle t accidentally similar)
\item
  \textbf{Early stopping} (for iterative models)
\end{enumerate}

If train error is also high: underfitting (high bias) --- need more features, more complex model, less regularization."

\subsection{10. Senior-Level Discussion Points}\label{10-senior-level-discussion-points}

\subsubsection{10.1 When to Choose Each Model in Production}\label{101-when-to-choose-each-model-in-production}

\textbf{The Reproducibility vs. Accuracy Tradeoff:}
A deep neural net may get 0.5\% better AUC than XGBoost but requires 10x more engineering for monitoring, retraining, debugging, and serving. For production systems, the simpler model that the team can maintain is often better.

\textbf{Feature engineering vs. architecture:}
In the era of AutoML and feature stores, the bottleneck has shifted from model complexity to feature quality. XGBoost with well-engineered features consistently beats neural networks on tabular data. The top Kaggle insight: feature engineering \textgreater\textgreater{} model choice for tabular data.

\subsubsection{10.2 Label Leakage and Data Pipeline Issues}\label{102-label-leakage-and-data-pipeline-issues}

More important than model choice: ensuring training/test splits don\textquotesingle t leak future information. Time-based splits for time-series. Target encoding with cross-validation to prevent leakage. Calibration checks --- log-loss can be low but probabilities systematically wrong (requires Platt scaling or isotonic regression).

\subsubsection{10.3 Class Imbalance at Scale}\label{103-class-imbalance-at-scale}

At FAANG scale, class imbalance is extreme: CTR prediction has 0.1\% positive rate, fraud detection 0.01\%. Techniques:

\begin{itemize}
\tightlist
\item
  \textbf{Undersampling} with calibration correction (Facebook\textquotesingle s negative downsampling + log-odds calibration)
\item
  \textbf{Class weights} (\texttt{class\_weight=\textquotesingle{}balanced\textquotesingle{}} in sklearn)
\item
  \textbf{Threshold tuning} --- optimize at deployment time based on business costs
\item
  \textbf{SMOTE} for small datasets only --- generates synthetic samples
\item
  \textbf{Focal loss} (Lin et al., RetinaNet) --- down-weights easy negatives in loss function
\end{itemize}

\subsubsection{10.4 Gradient Boosting vs Neural Networks for Tabular Data}\label{104-gradient-boosting-vs-neural-networks-for-tabular-data}

\textbf{Current consensus (Grinsztajn et al., 2022; Shwartz-Ziv \& Armon, 2022):} Tree-based models still dominate purely tabular tasks. Reasons:

\begin{itemize}
\tightlist
\item
  Tabular data is often "irregular" --- features have different scales, types, and distributions; trees are invariant to monotone transformations
\item
  Trees handle feature interactions naturally; neural nets need architectural tricks
\item
  Trees don\textquotesingle t suffer from uninformative features the same way
\end{itemize}

\textbf{When neural nets win:} Very large datasets (millions+ rows where boosting becomes slow), high-cardinality categoricals with learned embeddings, multi-modal inputs (tabular + text/image), need for transfer learning.

\subsubsection{10.5 The Decomposition of Gradient Boosting Loss}\label{105-the-decomposition-of-gradient-boosting-loss}

The XGBoost objective uses a second-order Taylor expansion of the loss:

\begin{verbatim}
L(y, F+h) ≈ L(y, F) + g*h + (1/2)*H*h^2

g = dL/dF    (gradient)
H = d^2L/dF^2   (Hessian)
\end{verbatim}

This enables optimal closed-form leaf weights: \texttt{w\_j*\ =\ -G\_j\ /\ (H\_j\ +\ lambda)} where G\_j and H\_j are gradient/Hessian sums in leaf j. The Hessian captures curvature --- it\textquotesingle s why XGBoost converges faster than first-order GBM.

\subsubsection{10.6 Feature Importance Pitfalls}\label{106-feature-importance-pitfalls}

\textbf{MDI (Mean Decrease in Impurity):}

\begin{itemize}
\tightlist
\item
  Biased toward high-cardinality features (more possible splits = more likely to appear useful)
\item
  Correlated features split importance between them, potentially understating each
\item
  Use permutation importance or SHAP for production feature analysis
\end{itemize}

\textbf{SHAP (SHapley Additive exPlanations):}

\begin{itemize}
\tightlist
\item
  Based on game theory: feature importance = average marginal contribution across all feature coalitions
\item
  TreeSHAP: O(T * L * D) --- computationally feasible for trees
\item
  Consistent and locally accurate by construction
\item
  The gold standard for model explainability at FAANG
\end{itemize}

\subsubsection{10.7 Hyperparameter Tuning Strategies}\label{107-hyperparameter-tuning-strategies}

\begin{longtable}[]{@{}
  >{\raggedright\arraybackslash}p{(\columnwidth - 2\tabcolsep) * \real{0.3955}}
  >{\raggedright\arraybackslash}p{(\columnwidth - 2\tabcolsep) * \real{0.6045}}@{}}
\toprule\noalign{}
\begin{minipage}[b]{\linewidth}\raggedright
Method
\end{minipage} & \begin{minipage}[b]{\linewidth}\raggedright
When to Use
\end{minipage} \\
\midrule\noalign{}
\endhead
\bottomrule\noalign{}
\endlastfoot
\textbf{Grid Search} & Small search space, deterministic needs \\
\textbf{Random Search} & Practical default --- often finds good params faster than grid search (Bergstra \& Bengio) \\
\textbf{Bayesian Optimization (Optuna, Hyperopt)} & Expensive models, important to find optimum \\
\textbf{Early Stopping + CV} & Tree models --- add trees until validation error stops improving \\
\textbf{Population-Based Training} & Parallel GPU training at Google DeepMind scale \\
\end{longtable}

\subsubsection{10.8 Calibration}\label{108-calibration}

A model with good AUC can have terrible calibration. In ads, you need calibrated probabilities for expected value calculation: \texttt{E{[}value{]}\ =\ P(click)\ *\ bid}. If P(click) is systematically 2x too high, you overbid and lose money.

Calibration methods:

\begin{itemize}
\tightlist
\item
  \textbf{Platt scaling:} Fit sigmoid on top of model output
\item
  \textbf{Isotonic regression:} Non-parametric, more flexible
\item
  \textbf{Temperature scaling:} For neural networks
\end{itemize}

\textbf{Reliability diagram:} Plot predicted probability bins vs actual event rate. Perfect calibration = diagonal line.

\subsection{11. Typical Mistakes Candidates Make}\label{11-typical-mistakes-candidates-make}

\subsubsection{Conceptual Mistakes}\label{conceptual-mistakes}

\begin{enumerate}
\def\labelenumi{\arabic{enumi}.}
\item
  \textbf{"Higher k in k-NN reduces overfitting"} --- Partially right: higher k reduces variance, but increases bias. The optimal k balances both.
\item
  \textbf{"Decision trees are always interpretable"} --- Only shallow trees (depth 3-4). A depth-20 tree with 1M nodes is as uninterpretable as a neural network.
\item
  \textbf{"SVM kernel picks the right feature space automatically"} --- The kernel defines the feature space; choosing the wrong kernel (e.g., RBF when data is linearly separable) can hurt performance.
\item
  \textbf{"Gradient boosting doesn\textquotesingle t overfit"} --- It absolutely can. More trees + high learning rate = overfit. Always use early stopping on a validation set.
\item
  \textbf{"PCA removes unimportant features"} --- PCA removes LOW-VARIANCE directions. Low variance \(\neq\) unimportant. If an important feature has low variance (binary label), PCA may remove it.
\item
  \textbf{"Logistic regression assumes the output is normally distributed"} --- No, it assumes the \textbf{log-odds} are linear in features. The output is a Bernoulli random variable.
\item
  \textbf{"Random Forest is just bagging"} --- Random Forest adds feature subsampling at each split (not just bootstrap rows). This is the critical innovation that decorrelates trees further.
\end{enumerate}

\subsubsection{Implementation Mistakes}\label{implementation-mistakes}

\begin{enumerate}
\def\labelenumi{\arabic{enumi}.}
\setcounter{enumi}{7}
\item
  \textbf{Not standardizing features for distance-based models} --- k-NN, SVM, PCA, k-Means all need feature scaling. A feature with range {[}0,1000{]} will dominate one with range {[}0,1{]}.
\item
  \textbf{Using test set to select k in k-NN} --- Data leakage. Use nested cross-validation or a separate validation set.
\item
  \textbf{Target encoding without cross-validation} --- Encoding target mean per category on full training set leaks label information. Use K-fold target encoding.
\item
  \textbf{Not handling class imbalance} --- Using accuracy as metric on 99/1 split (always predict majority \(\rightarrow\) 99\% accuracy, useless model). Use AUC, F1, or precision-recall.
\item
  \textbf{Comparing models on training accuracy} --- Always compare on held-out validation/test data. Training accuracy is only useful for diagnosing underfitting.
\item
  \textbf{k-Means++ initialization vs. random} --- Default in sklearn is k-means++. Saying "k-Means is sensitive to initialization" without mentioning k-means++ shows lack of depth.
\end{enumerate}

\subsubsection{Interview-Specific Mistakes}\label{interview-specific-mistakes}

\begin{enumerate}
\def\labelenumi{\arabic{enumi}.}
\setcounter{enumi}{13}
\item
  \textbf{Not asking clarifying questions} --- In a system design or ML design interview, jumping to XGBoost without asking about data size, latency requirements, interpretability needs shows poor engineering judgment.
\item
  \textbf{Ignoring data size constraints} --- SVM is O(n\^{}3) --- saying "I\textquotesingle d use SVM" for a billion-row dataset is a red flag.
\item
  \textbf{Not connecting to bias-variance} --- When asked about regularization, hyperparameters, or overfitting, always frame the answer in terms of the bias-variance tradeoff.
\end{enumerate}

\subsection{12. How This Connects To Other Topics}\label{12-how-this-connects-to-other-topics}

\subsubsection{ML Fundamentals}\label{ml-fundamentals}

\begin{itemize}
\tightlist
\item
  \textbf{Bias-Variance Tradeoff}: Trees (high variance) \(\rightarrow\) Random Forest (low variance) \(\rightarrow\) Understanding this is the foundation for hyperparameter tuning
\item
  \textbf{Optimization}: GD for logistic regression \(\rightarrow\) Adam for neural nets \(\rightarrow\) FTRL for online learning
\item
  \textbf{Feature engineering}: Classic models expose which features matter \(\rightarrow\) informs deep learning feature design
\item
  \textbf{Cross-validation}: Model selection for ALL algorithms; k-fold, stratified, time-series CV
\end{itemize}

\subsubsection{Deep Learning}\label{deep-learning}

\begin{itemize}
\tightlist
\item
  \textbf{Logistic Regression \(\rightarrow\) Neural Nets}: Logistic regression = neural net with no hidden layers. Adding hidden layers = deep learning.
\item
  \textbf{Gradient Boosting \(\rightarrow\) DNNs}: GBDT-NN hybrid (Meta\textquotesingle s approach) bridges the gap. DNNs learn representations; GBDTs learn feature interactions on tabular data.
\item
  \textbf{Kernels \(\rightarrow\) Radial Basis Function Networks}: RBF kernel SVMs are equivalent to shallow RBF networks.
\item
  \textbf{PCA \(\rightarrow\) Autoencoders}: Linear PCA is the special case of a linear autoencoder with identity activation. Non-linear autoencoders = kernel PCA generalization.
\item
  \textbf{Embeddings + k-NN}: Neural nets generate embeddings, k-NN (with ANN indexing) retrieves nearest neighbors --- this is the two-tower model architecture used in production recommendations.
\end{itemize}

\subsubsection{System Design}\label{system-design}

\begin{itemize}
\tightlist
\item
  \textbf{Inference latency}: Linear model: sub-millisecond. k-NN (naive): O(n\emph{p) --- slow for large n. FAISS approximate k-NN: O(log n). XGBoost: O(M}depth) --- typically fast.
\item
  \textbf{Training throughput}: Random Forest is embarrassingly parallel. Boosting is sequential. This affects distributed training architecture choices.
\item
  \textbf{Model size}: SVM stores support vectors (can be large); RF stores all trees (large); Logistic Regression = one weight vector (tiny).
\item
  \textbf{Online learning}: Logistic Regression with SGD supports online updates. Trees and SVMs do not (need retraining). Important for streaming ML systems.
\item
  \textbf{Feature stores}: Classic models require explicit feature engineering \(\rightarrow\) motivates centralized feature computation and caching infrastructure.
\item
  \textbf{A/B testing}: Statistical models (logistic regression) give p-values; gradient boosting requires bootstrap confidence intervals --- affects significance testing in experiments.
\end{itemize}

\subsection{13. FAANG Interview Tips}\label{13-faang-interview-tips}

\subsubsection{The ML Design Interview Framework}\label{the-ml-design-interview-framework}

When asked "design an ML system for X":

\begin{enumerate}
\def\labelenumi{\arabic{enumi}.}
\tightlist
\item
  \textbf{Clarify} --- What\textquotesingle s the objective metric? What\textquotesingle s the data size? Latency requirements? Interpretability needed?
\item
  \textbf{Baseline} --- Always start with Logistic Regression or a simple GBM as baseline
\item
  \textbf{Features} --- Spend most time here; features matter more than model choice
\item
  \textbf{Model selection} --- Justify based on data size, latency, interpretability, expected non-linearity
\item
  \textbf{Evaluation} --- Offline metrics (AUC, NDCG) vs. online metrics (CTR, conversion)
\item
  \textbf{Production concerns} --- Distribution shift, retraining frequency, monitoring
\end{enumerate}

\subsubsection{Coding Interviews}\label{coding-interviews}

You may be asked to implement from scratch:

\begin{Shaded}
\begin{Highlighting}[]
\CommentTok{\# Linear Regression (gradient descent)}
\KeywordTok{def}\NormalTok{ fit(X, y, lr}\OperatorTok{=}\FloatTok{0.01}\NormalTok{, epochs}\OperatorTok{=}\DecValTok{1000}\NormalTok{):}
\NormalTok{    n, p }\OperatorTok{=}\NormalTok{ X.shape}
\NormalTok{    w }\OperatorTok{=}\NormalTok{ np.zeros(p)}
    \ControlFlowTok{for}\NormalTok{ \_ }\KeywordTok{in} \BuiltInTok{range}\NormalTok{(epochs):}
\NormalTok{        y\_hat }\OperatorTok{=}\NormalTok{ X }\OperatorTok{@}\NormalTok{ w}
\NormalTok{        grad }\OperatorTok{=}\NormalTok{ (}\DecValTok{2}\OperatorTok{/}\NormalTok{n) }\OperatorTok{*}\NormalTok{ X.T }\OperatorTok{@}\NormalTok{ (y\_hat }\OperatorTok{{-}}\NormalTok{ y)}
\NormalTok{        w }\OperatorTok{{-}=}\NormalTok{ lr }\OperatorTok{*}\NormalTok{ grad}
    \ControlFlowTok{return}\NormalTok{ w}

\CommentTok{\# k{-}NN Classification}
\KeywordTok{def}\NormalTok{ predict(X\_train, y\_train, x\_test, k}\OperatorTok{=}\DecValTok{5}\NormalTok{):}
\NormalTok{    distances }\OperatorTok{=}\NormalTok{ np.sqrt(np.}\BuiltInTok{sum}\NormalTok{((X\_train }\OperatorTok{{-}}\NormalTok{ x\_test)}\OperatorTok{**}\DecValTok{2}\NormalTok{, axis}\OperatorTok{=}\DecValTok{1}\NormalTok{))}
\NormalTok{    k\_indices }\OperatorTok{=}\NormalTok{ np.argsort(distances)[:k]}
\NormalTok{    k\_labels }\OperatorTok{=}\NormalTok{ y\_train[k\_indices]}
    \ControlFlowTok{return}\NormalTok{ np.bincount(k\_labels).argmax()}

\CommentTok{\# k{-}Means}
\KeywordTok{def}\NormalTok{ kmeans(X, k, max\_iter}\OperatorTok{=}\DecValTok{100}\NormalTok{):}
\NormalTok{    centroids }\OperatorTok{=}\NormalTok{ X[np.random.choice(}\BuiltInTok{len}\NormalTok{(X), k, replace}\OperatorTok{=}\VariableTok{False}\NormalTok{)]}
    \ControlFlowTok{for}\NormalTok{ \_ }\KeywordTok{in} \BuiltInTok{range}\NormalTok{(max\_iter):}
        \CommentTok{\# Assignment}
\NormalTok{        dists }\OperatorTok{=}\NormalTok{ np.sqrt(((X }\OperatorTok{{-}}\NormalTok{ centroids[:, np.newaxis])}\OperatorTok{**}\DecValTok{2}\NormalTok{).}\BuiltInTok{sum}\NormalTok{(axis}\OperatorTok{=}\DecValTok{2}\NormalTok{))}
\NormalTok{        labels }\OperatorTok{=}\NormalTok{ np.argmin(dists, axis}\OperatorTok{=}\DecValTok{0}\NormalTok{)}
        \CommentTok{\# Update}
\NormalTok{        new\_centroids }\OperatorTok{=}\NormalTok{ np.array([X[labels}\OperatorTok{==}\NormalTok{i].mean(}\DecValTok{0}\NormalTok{) }\ControlFlowTok{for}\NormalTok{ i }\KeywordTok{in} \BuiltInTok{range}\NormalTok{(k)])}
        \ControlFlowTok{if}\NormalTok{ np.allclose(centroids, new\_centroids):}
            \ControlFlowTok{break}
\NormalTok{        centroids }\OperatorTok{=}\NormalTok{ new\_centroids}
    \ControlFlowTok{return}\NormalTok{ labels, centroids}

\CommentTok{\# Logistic Regression}
\KeywordTok{def}\NormalTok{ sigmoid(z):}
    \ControlFlowTok{return} \DecValTok{1} \OperatorTok{/}\NormalTok{ (}\DecValTok{1} \OperatorTok{+}\NormalTok{ np.exp(}\OperatorTok{{-}}\NormalTok{z))}

\KeywordTok{def}\NormalTok{ fit\_logistic(X, y, lr}\OperatorTok{=}\FloatTok{0.1}\NormalTok{, epochs}\OperatorTok{=}\DecValTok{1000}\NormalTok{):}
\NormalTok{    n, p }\OperatorTok{=}\NormalTok{ X.shape}
\NormalTok{    w }\OperatorTok{=}\NormalTok{ np.zeros(p)}
    \ControlFlowTok{for}\NormalTok{ \_ }\KeywordTok{in} \BuiltInTok{range}\NormalTok{(epochs):}
\NormalTok{        p\_hat }\OperatorTok{=}\NormalTok{ sigmoid(X }\OperatorTok{@}\NormalTok{ w)}
\NormalTok{        grad }\OperatorTok{=}\NormalTok{ (}\DecValTok{1}\OperatorTok{/}\NormalTok{n) }\OperatorTok{*}\NormalTok{ X.T }\OperatorTok{@}\NormalTok{ (p\_hat }\OperatorTok{{-}}\NormalTok{ y)}
\NormalTok{        w }\OperatorTok{{-}=}\NormalTok{ lr }\OperatorTok{*}\NormalTok{ grad}
    \ControlFlowTok{return}\NormalTok{ w}
\end{Highlighting}
\end{Shaded}

\subsubsection{Typical FAANG Question Flows}\label{typical-faang-question-flows}

\textbf{"Which model would you use for CTR prediction at ads scale?"}
\(\rightarrow\) "Start with Logistic Regression for speed and interpretability (FTRL for online learning). Move to GBDT for capturing feature interactions (XGBoost or LightGBM). At scale, use a hybrid: GBDT to generate crossed features fed into a neural net (like Meta\textquotesingle s DLRM). Key considerations: training data is billions of rows, serving latency \textless{} 5ms, need calibrated probabilities for expected value computation."

\textbf{"How does XGBoost handle missing values?"}
\(\rightarrow\) "During training, XGBoost learns a default direction for missing values --- for each split, it tries routing missing values to left vs. right and picks whichever reduces the loss more. At inference, missing values follow the learned default direction. This is learned automatically, not imputed."

\textbf{"Why does bagging work but independent models don\textquotesingle t?"}
\(\rightarrow\) "Independence is the key. Training independent models on the same dataset gives correlated predictions --- their errors don\textquotesingle t cancel. Bootstrap sampling + random feature subsampling reduces the correlation between trees\textquotesingle{} errors. The variance formula shows this clearly: \texttt{Var(avg)\ =\ rho\ *\ sigma\^{}2\ +\ (1-rho)/B\ *\ sigma\^{}2}. Independent models have \texttt{rho=0} but you\textquotesingle d need separate data. Bootstrap gives \texttt{rho\textgreater{}0} but real data, which is the practical trade-off."

\subsection{14. Revision Cheat Sheet}\label{14-revision-cheat-sheet}

\subsubsection{10-Minute Summary}\label{10-minute-summary}

\textbf{Linear Regression:} Fit line by minimizing squared errors. OLS = \texttt{w\ =\ (X\^{}T\ X)\^{}\{-1\}\ X\^{}T\ y}. Ridge: L2 reg, no zeros. Lasso: L1 reg, sparsity. Assumptions: linearity, no multicollinearity, homoscedasticity.

\textbf{Logistic Regression:} \texttt{P(y=1\textbar{}x)\ =\ sigmoid(w\^{}T\ x)}. Loss = log-loss (cross-entropy). No closed form. Linear boundary. Coefficients = log-odds change per unit feature.

\textbf{Decision Tree:} Recursively split on best feature (Gini or IG). Greedy, high variance, interpretable. Prune with max\_depth, min\_samples\_leaf. CART algorithm.

\textbf{Random Forest:} B trees on bootstrap samples + random feature subset at each split. Reduces variance without increasing bias. OOB error = free validation. Feature importance via MDI or permutation.

\textbf{Gradient Boosting:} Sequential trees fitting pseudo-residuals (negative gradient). XGBoost adds L2 tree regularization + second-order Taylor + sparse-aware. LightGBM adds leaf-wise growth + histograms + GOSS + EFB. Fastest on tabular data.

\textbf{SVM:} Maximize margin hyperplane. Support vectors define it. Soft margin: C trades margin width vs. violations. Kernel trick: implicit high-dim mapping via inner products. RBF = infinite-dim, O(p) computation.

\textbf{k-NN:} Store all training data. Predict by k-nearest votes/mean. Lazy learning. Curse of dimensionality. Need feature scaling. k-means++ init. Use ANN (FAISS) for scale.

\textbf{Naive Bayes:} \texttt{P(y\textbar{}x)\ ∝\ P(y)\ *\ prod\ P(x\_j\textbar{}y)}. Assumes conditional independence. Fast, works on text. Gaussian/Multinomial/Bernoulli variants. Laplace smoothing for zero probabilities.

\textbf{k-Means:} Assign points to nearest centroid, update centroids as mean. WCSS objective, non-convex. k-means++ init. Elbow/silhouette for choosing k. Spherical clusters assumption.

\textbf{Hierarchical:} Build dendrogram bottom-up (agglomerative). Ward linkage = minimize variance increase. Cut dendrogram to get k clusters. O(n\^{}2) memory.

\textbf{DBSCAN:} Core/border/noise points. Finds arbitrary shapes. No need to specify k. Sensitive to eps and min\_samples. Natural outlier detection.

\textbf{PCA:} Eigenvectors of covariance matrix = principal components. Project onto top-k for dimensionality reduction. Use SVD in practice. Standardize first. Explained variance ratio for choosing k.

\textbf{Ensembles:} Bagging = parallel, reduces variance. Boosting = sequential, reduces bias. Stacking = meta-model on predictions.

\subsubsection{Algorithm Decision Table --- "Which Algorithm When"}\label{algorithm-decision-table--which-algorithm-when}

\begin{longtable}[]{@{}
  >{\raggedright\arraybackslash}p{(\columnwidth - 4\tabcolsep) * \real{0.3175}}
  >{\raggedright\arraybackslash}p{(\columnwidth - 4\tabcolsep) * \real{0.3209}}
  >{\raggedright\arraybackslash}p{(\columnwidth - 4\tabcolsep) * \real{0.3615}}@{}}
\toprule\noalign{}
\begin{minipage}[b]{\linewidth}\raggedright
Scenario
\end{minipage} & \begin{minipage}[b]{\linewidth}\raggedright
Recommended
\end{minipage} & \begin{minipage}[b]{\linewidth}\raggedright
Why
\end{minipage} \\
\midrule\noalign{}
\endhead
\bottomrule\noalign{}
\endlastfoot
\textbf{Tabular data, maximize accuracy} & LightGBM / XGBoost & State-of-the-art on tabular; handles all feature types \\
\textbf{Need interpretability / regulatory compliance} & Logistic Regression or Decision Tree (shallow) & Explicit feature coefficients or printable rules \\
\textbf{Baseline model} & Logistic Regression & Fast, calibrated, easy to improve upon \\
\textbf{High-dimensional sparse features (text/NLP)} & Logistic Regression (L1/L2) or Linear SVM & Efficient on sparse data; SVM with linear kernel is L2-LR \\
\textbf{Small dataset (n \textless{} 1000), many features} & SVM (RBF kernel) & Effective in high p/n ratio; global optimum \\
\textbf{Non-linear patterns, medium dataset} & SVM + RBF kernel or Random Forest & Kernel SVM for structured non-linearity; RF for robustness \\
\textbf{Very large dataset (n \textgreater{} 1M), need speed} & LightGBM & Fastest tree-based; designed for large-scale \\
\textbf{Text classification, real-time} & Naive Bayes & Sub-millisecond inference; surprisingly competitive \\
\textbf{Need calibrated probabilities} & Logistic Regression or calibrated GBM & LR is naturally calibrated; GBM + Platt scaling \\
\textbf{Missing values, heterogeneous features} & XGBoost / LightGBM & Native missing value handling; no preprocessing \\
\textbf{Unknown number of clusters, outlier detection} & DBSCAN & Finds k automatically; labels noise points \\
\textbf{Known k, spherical clusters} & k-Means & Fast, interpretable, works well with good init \\
\textbf{Cluster hierarchy / exploratory analysis} & Hierarchical (Ward linkage) & Dendrogram shows structure at all levels \\
\textbf{Nearest neighbor retrieval at scale} & FAISS / HNSW (ANN) & Exact k-NN is O(n*p); ANN is O(log n) \\
\textbf{Recommendation (item-item, user-item)} & k-NN in embedding space & After generating embeddings with Matrix Factorization/DNN \\
\textbf{Remove multicollinearity before linear model} & PCA & Creates orthogonal features \\
\textbf{Visualization of high-dim data} & t-SNE / UMAP & Preserves local/global structure; PCA first for large n \\
\textbf{Feature compression (autoencoder alternative)} & PCA & Linear compression; interpretable explained variance \\
\textbf{Online/streaming learning} & Logistic Regression (SGD) & Tree models can\textquotesingle t update incrementally \\
\textbf{Maximum accuracy, no latency constraint} & Stacking (GBM + RF + LR meta) & Combines diverse model predictions \\
\textbf{Interpretable feature interactions} & Decision Tree \(\rightarrow\) tree.plot() & Explicit IF-THEN rules for business users \\
\textbf{Imbalanced classes, small minority} & XGBoost + class\_weight + threshold tuning & Robust to imbalance; easy threshold calibration \\
\textbf{Anomaly detection (labeled)} & Logistic Regression / Isolation Forest & LR if labeled; IF for unsupervised \\
\textbf{Anomaly detection (unlabeled)} & DBSCAN / Isolation Forest / Autoencoder & DBSCAN: density-based; IF: tree-based isolation \\
\end{longtable}

\subsubsection{Complexity Quick Reference}\label{complexity-quick-reference}

\begin{longtable}[]{@{}
  >{\raggedright\arraybackslash}p{(\columnwidth - 6\tabcolsep) * \real{0.2568}}
  >{\raggedright\arraybackslash}p{(\columnwidth - 6\tabcolsep) * \real{0.2838}}
  >{\raggedright\arraybackslash}p{(\columnwidth - 6\tabcolsep) * \real{0.3514}}
  >{\raggedright\arraybackslash}p{(\columnwidth - 6\tabcolsep) * \real{0.1081}}@{}}
\toprule\noalign{}
\begin{minipage}[b]{\linewidth}\raggedright
Algorithm
\end{minipage} & \begin{minipage}[b]{\linewidth}\raggedright
Train Time
\end{minipage} & \begin{minipage}[b]{\linewidth}\raggedright
Inference Time
\end{minipage} & \begin{minipage}[b]{\linewidth}\raggedright
Memory
\end{minipage} \\
\midrule\noalign{}
\endhead
\bottomrule\noalign{}
\endlastfoot
Linear Regression & O(n*p\^{}2) closed form & O(p) & O(p) \\
Logistic Regression & O(n\emph{p}iter) & O(p) & O(p) \\
Decision Tree & O(n\emph{p}log\^{}2 n) & O(depth) & O(n) \\
Random Forest & O(B\emph{n}sqrt(p)*log\^{}2 n) & O(B*depth) & O(B*n) \\
XGBoost & O(M\emph{n}p*log n) & O(M*depth) & O(n*p) \\
LightGBM & O(M\emph{n}255*log n) & O(M*depth) & O(n*255) \\
SVM & O(n\^{}2 p to n\^{}3 p) & O(n\_sv * p) & O(n\_sv) \\
k-NN & O(n*p) store & O(n*p + n log n) per query & O(n*p) \\
Naive Bayes & O(n*p) & O(p*K) & O(p*K) \\
k-Means & O(n\emph{k}p*iter) & O(k*p) & O(n*k) \\
Hierarchical & O(n\^{}2 log n) & O(1) post & O(n\^{}2) \\
DBSCAN & O(n log n) with index & O(1) & O(n) \\
PCA & O(n\emph{p}k) randomized & O(p*k) & O(p*k) \\
\end{longtable}

\subsubsection{Key Formulas Cheat Sheet}\label{key-formulas-cheat-sheet}

\begin{verbatim}
LOSS FUNCTIONS:
MSE (regression):     L = (1/n) * sum[(y_i - y_hat_i)^2]
Log-loss (binary):    L = -(1/n) * sum[y*log(p) + (1-y)*log(1-p)]
Hinge (SVM):          L = (1/n) * sum[max(0, 1 - y_i * f(x_i))]

IMPURITY MEASURES:
Gini:     G = 1 - sum_k[p_k^2]          (0 = pure, 0.5 = worst binary)
Entropy:  H = -sum_k[p_k * log2(p_k)]   (0 = pure, 1 = worst binary)

REGULARIZATION:
Ridge: L = MSE + lambda * sum(w_j^2)
Lasso: L = MSE + lambda * sum(|w_j|)
SVM:   min (1/2)||w||^2 + C * sum[xi_i]

PROBABILITY:
Sigmoid:  sigma(z) = 1/(1 + e^{-z})
Softmax:  P(y=k) = e^{w_k^T x} / sum_j[e^{w_j^T x}]
Bayes:    P(y|x) ∝ P(x|y) * P(y)

KERNELS:
Linear:   K(x,z) = x^T z
RBF:      K(x,z) = exp(-gamma * ||x-z||^2)
Poly:     K(x,z) = (gamma * x^T z + r)^d

DIMENSIONALITY:
PCA:      Sigma v = lambda v    (eigenvectors of covariance)
EVR:      EVR_k = lambda_k / sum_i lambda_i

VARIANCE OF ENSEMBLE:
Var(avg B corr trees) = rho*sigma^2 + (1-rho)*sigma^2/B
\end{verbatim}

\subsubsection{Most Important Concepts --- Final 5}\label{most-important-concepts--final-5}

\begin{enumerate}
\def\labelenumi{\arabic{enumi}.}
\item
  \textbf{Gradient boosting = gradient descent in function space} --- Each tree corrects the residual error of the ensemble so far; sequential trees are steps along the loss gradient.
\item
  \textbf{Kernel trick = inner product without feature expansion} --- \texttt{K(x,z)\ =\ \textless{}phi(x),\ phi(z)\textgreater{}} computed in O(p) without computing phi, enabling infinite-dimensional feature spaces.
\item
  \textbf{Bagging reduces variance by decorrelating trees; boosting reduces bias by sequential correction} --- Know these two error types, know which ensemble addresses each.
\item
  \textbf{Bias-Variance Decomposition frames everything} --- Every hyperparameter choice (tree depth, k in k-NN, C in SVM, lambda in regularization) is trading off bias vs. variance. Always frame your answers through this lens.
\item
  \textbf{The "right algorithm" depends on constraints} --- Data size (n, p), latency budget, interpretability requirements, missing values, class balance. There is no universally best algorithm --- demonstrating this judgment is what FAANG interviewers want.
\end{enumerate}

\emph{Word count: \textasciitilde9,800 words \textbar{} Section count: 14 (Overview through Revision Cheat Sheet)}

\clearpage
\section{Visual Reference}
\noindent A set of figures distilling the key quantitative intuitions of this topic.\par\vspace{4pt}
\visualfigure{../figures/aiml-02-classic/01-roc-pr.pdf}{ROC summarizes ranking quality across thresholds; PR is more informative under heavy class imbalance, where a high AUC can still hide poor precision.}
\end{document}
